%============================================================
%  Bd. 11 - Äquivalenzrelationen
%============================================================

\documentclass[main.tex]{subfiles}

\ifSubfilesClassLoaded{
  \BandSetupStandalone{B11}
}{
}

\title{Bd. 11 - Äquivalenzrelationen}
\author{Martin Kunze}
\date{}
\setcounter{file}{11}

\hypersetup{
  pdftitle={Bd. 11 - Äquivalenzrelationen},
  pdfauthor={Martin Kunze}
}

\begin{document}

\title{Bd. 11 - Äquivalenzrelationen}
\author{Martin Kunze}
\date{}
\hypersetup{pageanchor=false}
\maketitle
\hypersetup{pageanchor=true}
\setcounter{tocdepth}{3}
\tableofcontents
\setcounter{file}{11}
\renewcommand{\FormulaBandID}{11}

\chapter{Einleitung}

Äquivalenzrelationen strukturieren einen Träger durch Klassen
ununterscheidbarer Elemente. Nach dem allgemeinen Begriff werden in einem
eigenen Kapitel die Ununterscheidbarkeitsrelation und die Gleichmächtigkeit
als Beispiele vollständig aufgebaut. Darauf folgen Äquivalenzklassen und
Quotientenmengen. Das abschließende Kapitel entwickelt die punktweise
Quotientenprojektion und die von ihr induzierte Quotientenbildabbildung; der
Ununterscheidbarkeitsquotient spezialisiert beide Konstruktionen.

\chapter{Begriff der Äquivalenzrelation}

\section{Axiome einer Äquivalenzrelation}

\FormulaAxiomDeltaK[Reflexivität]{%
  x\in A \vdash x\sim x%
}{EqRelReflexivity}{
  \DeltaRow{Mengen}{A \dsep x}
  \DeltaRow{Relationsoperatoren}{\sim}
}

\FormulaAxiomDeltaK[Symmetrie]{%
  x\sim y \vdash y\sim x%
}{EqRelSymmetry}{
  \DeltaRow{Mengen}{A \dsep x \dsep y}
  \DeltaRow{Relationsoperatoren}{\sim}
}

\FormulaAxiomDeltaK[Transitivität]{%
  x\sim y \dsep y\sim z \vdash x\sim z%
}{EqRelTransitivity}{
  \DeltaRow{Mengen}{A \dsep x \dsep y \dsep z}
  \DeltaRow{Relationsoperatoren}{\sim}
}

\section{Definition der Äquivalenzrelation}

\FormulaDefDeltaK[Begriff der Äquivalenzrelation]{\EqRel{A,\sim}}{EqRelDef}{
  \DeltaRow{Mengen}{A \dsep x \dsep y \dsep z}
  \DeltaRow{Relationsoperatoren}{\sim}
  \DeltaRow{\textbf{Axiome}}{}
  \DeltaRow{Binärer Relationsoperator}
           {\sim\ \text{auf }A}
  \DeltaRow{Reflexivität}
           {x\in A \vdash x\sim x}
           [\FormulaRefAuto{EqRelReflexivity}[ax]]
  \DeltaRow{Symmetrie}
           {x\sim y \vdash y\sim x}
           [\FormulaRefAuto{EqRelSymmetry}[ax]]
  \DeltaRow{Transitivität}
           {x\sim y \dsep y\sim z \vdash x\sim z}
           [\FormulaRefAuto{EqRelTransitivity}[ax]]
  \DeltaRow{\textbf{Neue Symbole}}{}
  \DeltaPrem{Äquivalenzrelation}{\EqRel{A,\sim}}
}

Eine Äquivalenzrelation wird damit nicht durch einen austauschbaren
Buchstaben \(R\), sondern durch ihren binären Relationsoperator bezeichnet.
Die Aussage \(x\sim y\) ist seine durchgängig verwendete Infixnotation.

\chapter{Beispiele für Äquivalenzrelationen}

\section{Ununterscheidbarkeitsrelation einer Mengenfamilie}

Ein wichtiges Beispiel entsteht aus einer Mengenfamilie \(M\). In der
Terminologie der Rough-Set-Theorie ist dies eine
Ununterscheidbarkeitsrelation (indiscernibility relation): Zwei Elemente des
Trägers \(\bigcup M\) sind bezüglich \(M\) ununterscheidbar, wenn sie in jedem
Mitglied von \(M\) entweder beide enthalten oder beide nicht enthalten sind.

\FormulaDefDeltaK[Ununterscheidbarkeitsoperator einer Mengenfamilie]{%
  \begin{aligned}[t]
    &a,b\in\bigcup M\vdash{}\\
    &a\mathrel{\SameOccRel{M}}b
      \coloneqq \forall X\in M\,(a\in X\leftrightarrow b\in X)
  \end{aligned}
}{SameOccRelDef}{%
  \DeltaRow{Mengen}{M\dsep X\dsep a\dsep b}
  \DeltaRow{Relationsoperatoren}{\SameOccRel{M}}
  \DeltaRow{Neue Symbole}{\SameOccRel{M}}
}

\FormulaThmDeltaK[Charakterisierung des Ununterscheidbarkeitsoperators]{%
  a,b\in\bigcup M\vdash
  \Bigl(a\mathrel{\SameOccRel{M}}b
  \leftrightarrow
  \forall X\in M\,(a\in X\leftrightarrow b\in X)\Bigr)
}{SameOccRelChar}{%
  \DeltaRow{Mengen}{M\dsep X\dsep a\dsep b}
  \DeltaRow{Relationsoperatoren}{\SameOccRel{M}}
}
\begin{tabproofwide}
  \proofstepwidestar[1]{a\in\bigcup M}{\rA}
  \proofstepwidestar[2]{b\in\bigcup M}{\rA}
  \proofstepwidestar[1,2]{%
    a\mathrel{\SameOccRel{M}}b
    \leftrightarrow
    \forall X\in M\,(a\in X\leftrightarrow b\in X)}{%
    \FormulaRefAuto{SameOccRelDef}[def]{1,2}}
\end{tabproofwide}

\FormulaThmDeltaK[Reflexivität der Ununterscheidbarkeitsrelation]{%
  a\in\bigcup M \vdash a\mathrel{\SameOccRel{M}}a
}{SameOccRelReflexivity}{%
  \DeltaRow{Mengen}{M\dsep X\dsep a}
  \DeltaRow{Relationsoperatoren}{\SameOccRel{M}}
}
\begin{tabproof}
  \proofstep{1}{a\in\bigcup M}{\rA}
  \proofstep{2}{X\in M}{\rA}
  \proofstep{1,2}{a\in X\leftrightarrow a\in X}{\FormulaRefAuto{P\leftrightarrow P}}
  \proofstep{1}{\forall X\in M\,(a\in X\leftrightarrow a\in X)}{\rUI{\rRI{2,3}}}
  \proofstep{1}{a\mathrel{\SameOccRel{M}}a}{%
    \FormulaRefAuto{P\leftrightarrow Q, Q\vdash P}{%
      \FormulaRefAuto{SameOccRelChar}[thm]{1,1},4}}
\end{tabproof}

\FormulaThmDeltaK[Symmetrie der Ununterscheidbarkeitsrelation]{%
  a\mathrel{\SameOccRel{M}}b \vdash b\mathrel{\SameOccRel{M}}a
}{SameOccRelSymmetry}{%
  \DeltaRow{Mengen}{M\dsep X\dsep a\dsep b}
  \DeltaRow{Relationsoperatoren}{\SameOccRel{M}}
}
\begin{tabproofwide}
  \proofstepwidestar[1]{a\mathrel{\SameOccRel{M}}b}{\rA}
  \proofstepwidestar[2]{a\in\bigcup M}{\rA}
  \proofstepwidestar[3]{b\in\bigcup M}{\rA}
  \proofstepwidestar[1,2,3]{%
    \forall X\in M\,(a\in X\leftrightarrow b\in X)}{%
    \FormulaRefAuto{P\leftrightarrow Q, P\vdash Q}{%
      \FormulaRefAuto{SameOccRelChar}[thm]{2,3},1}}
  \proofstepwidestar[1,2,3]{\forall X\in M\,(b\in X\leftrightarrow a\in X)}{%
    \FormulaRefAuto{\forall x(P(x)\leftrightarrow Q(x))\vdash
    \forall x(Q(x)\leftrightarrow P(x))}{4}}
  \proofstepwidestar[1,2,3]{b\mathrel{\SameOccRel{M}}a}{%
    \FormulaRefAuto{P\leftrightarrow Q, Q\vdash P}{%
      \FormulaRefAuto{SameOccRelChar}[thm]{3,2},5}}
\end{tabproofwide}

\FormulaThmDeltaK[Transitivität der Ununterscheidbarkeitsrelation]{%
  a\mathrel{\SameOccRel{M}}b\dsep b\mathrel{\SameOccRel{M}}c
  \vdash a\mathrel{\SameOccRel{M}}c
}{SameOccRelTransitivity}{%
  \DeltaRow{Mengen}{M\dsep X\dsep a\dsep b\dsep c}
  \DeltaRow{Relationsoperatoren}{\SameOccRel{M}}
}
\begin{tabproofwide}
  \proofstepwidestar[1]{a\mathrel{\SameOccRel{M}}b}{\rA}
  \proofstepwidestar[2]{b\mathrel{\SameOccRel{M}}c}{\rA}
  \proofstepwidestar[3]{a\in\bigcup M}{\rA}
  \proofstepwidestar[4]{b\in\bigcup M}{\rA}
  \proofstepwidestar[5]{c\in\bigcup M}{\rA}
  \proofstepwidestar[1,3,4]{%
    \forall X\in M\,(a\in X\leftrightarrow b\in X)}{%
    \FormulaRefAuto{P\leftrightarrow Q, P\vdash Q}{%
      \FormulaRefAuto{SameOccRelChar}[thm]{3,4},1}}
  \proofstepwidestar[2,4,5]{%
    \forall X\in M\,(b\in X\leftrightarrow c\in X)}{%
    \FormulaRefAuto{P\leftrightarrow Q, P\vdash Q}{%
      \FormulaRefAuto{SameOccRelChar}[thm]{4,5},2}}
  \proofstepwidestar[1,2,3,4,5]{%
    \forall X\in M\,(a\in X\leftrightarrow c\in X)}{%
    \FormulaRefAuto{\forall x(P(x)\leftrightarrow Q(x)),
    \forall x(Q(x)\leftrightarrow R(x))\vdash
    \forall x(P(x)\leftrightarrow R(x))}{6,7}}
  \proofstepwidestar[1,2,3,4,5]{a\mathrel{\SameOccRel{M}}c}{%
    \FormulaRefAuto{P\leftrightarrow Q, Q\vdash P}{%
      \FormulaRefAuto{SameOccRelChar}[thm]{3,5},8}}
\end{tabproofwide}

\FormulaThmDeltaK[Ununterscheidbarkeitsrelation ist Äquivalenzrelation]{%
  \EqRel{\bigcup M,\SameOccRel{M}}
}{SameOccRelEqRel}{%
  \DeltaRow{Mengen}{M\dsep X\dsep a\dsep b\dsep c}
  \DeltaRow{Relationsoperatoren}{\SameOccRel{M}}
}
\begin{tabproof}
  \proofstep{}{a\in\bigcup M \rightarrow a\mathrel{\SameOccRel{M}}a}{%
    \FormulaRefAuto{SameOccRelReflexivity}}
  \proofstep{}{%
    a\mathrel{\SameOccRel{M}}b \rightarrow b\mathrel{\SameOccRel{M}}a}{%
    \FormulaRefAuto{SameOccRelSymmetry}}
  \proofstep{}{%
    \bigl(a\mathrel{\SameOccRel{M}}b\land
          b\mathrel{\SameOccRel{M}}c\bigr)
    \rightarrow a\mathrel{\SameOccRel{M}}c}{%
    \FormulaRefAuto{SameOccRelTransitivity}}
  \proofstep{}{\EqRel{\bigcup M,\SameOccRel{M}}}{%
    \FormulaRefAuto{EqRelDef}{1,2,3}}
\end{tabproof}

\section{Gleichmächtigkeit}

Gleichmächtigkeit wird hier nicht als neuer primitiver Begriff axiomatisiert,
sondern durch die Existenz einer Bijektion definiert. Ihre grundlegenden
Eigenschaften gehören deshalb in diesen Band: Sie zeigen unmittelbar, dass
Gleichmächtigkeit eine Äquivalenzrelation auf jeder Mengenfamilie ist. Die
funktionentheoretischen Aussagen über induzierte Potenzmengenabbildungen
bleiben in Band~08; im vorliegenden Abschnitt werden daraus nur die
Gleichmächtigkeitsfolgen gewonnen.

\FormulaDefDeltaK[Begriff der gleichmächtigen Mengen]{A \EqCard B \coloneqq \exists F\colon A\bij B}{Gleichmächtigkeit}{
  \DeltaRow{Mengen}{A\dsep B}
  \DeltaRow{Funktionen}{F}
  \DeltaRow{\textbf{Neue Symbole}}{}
  \DeltaRow{Gleichmächtigkeit}{\EqCard}
}

\subsection{Grundlegende Eigenschaften}

\FormulaThmDeltaK[Reflexivität der Gleichmächtigkeit]{%
  A \EqCard A%
}{EqCardReflexive}{
  \DeltaRow{Mengen}{A}
}
\begin{tabproofsplit}
  \proofpartindR[Identitätsbijektion als Zeuge]{%
    \exists F\colon A\bij A%
  }{%
    \exists F\colon A\bij A%
  }[EqCardReflexiveBijectionWitness]
    \proofstep{}{\Id_A\colon A\bij A}{%
      \FormulaRefAuto{\Id_A\colon A\bij A}}
    \proofstep{}{\exists F\colon A\bij A}{\rEI{1}}
  \closeproofpartindR

  \proofpartindR[Übergang zur Gleichmächtigkeit]{%
    A\EqCard A%
  }{%
    A\EqCard A%
  }[EqCardReflexiveConclusion]
    \proofstep{}{\exists F\colon A\bij A}{%
      \FormulaRefAuto{EqCardReflexiveBijectionWitness}}
    \proofstep{}{A\EqCard A}{%
      \FormulaRefAuto{A\EqCard B\coloneqq
        \exists F\colon A\bij B}{1}}
  \closeproofpartindR
\end{tabproofsplit}

\FormulaThmDeltaK[Symmetrie der Gleichmächtigkeit]{%
  A \EqCard B \vdash B \EqCard A%
}{EqCardSymmetric}{
  \DeltaRow{Mengen}{A\dsep B}
  \DeltaRow{Funktionen}{F\dsep G}
}
\begin{tabproof}
  \proofstep{1}{A \EqCard B}{\rA}
  \proofstep{1}{\exists F\colon A\bij B}{\FormulaRefAuto{A \EqCard B \coloneqq \exists F\colon A\bij B}{1}}

  \proofstep{3}{F\colon A\bij B}{\rA}
  \proofstep{3}{F^{-1}\colon B\bij A}{\FormulaRefAuto{F\colon A\bij B\vdash F^{-1}\colon B\bij A}{3}}
  \proofstep{3}{\exists G\colon B\bij A}{\rEI{4}}

  \proofstep{1}{\exists G\colon B\bij A}{\rEE{2,3,5}}
  \proofstep{1}{B \EqCard A}{\FormulaRefAuto{A \EqCard B \coloneqq \exists F\colon A\bij B}{6}}
\end{tabproof}

\FormulaThmDeltaK[Transitivität der Gleichmächtigkeit]{%
  A \EqCard B \dsep B \EqCard C \vdash A \EqCard C%
}{EqCardTransitive}{
  \DeltaRow{Mengen}{A\dsep B\dsep C}
  \DeltaRow{Funktionen}{F\dsep G\dsep H}
}
\begin{tabproofsplit}
  \proofpartindR[Komposition fester Bijektionen]{%
    F\colon A\bij B\dsep G\colon B\bij C
    \vdash \exists H\colon A\bij C%
  }{%
    F\colon A\bij B\dsep G\colon B\bij C
    \vdash \exists H\colon A\bij C%
  }[EqCardTransitiveFixedBijections]
    \proofstep{1}{F\colon A\bij B}{\rA}
    \proofstep{2}{G\colon B\bij C}{\rA}
    \proofstep{1,2}{G\circ F\colon A\bij C}{%
      \FormulaRefAuto{F\colon A\bij B\dsep G\colon B\bij C
        \vdash G\circ F\colon A\bij C}{1,2}}
    \proofstep{1,2}{\exists H\colon A\bij C}{\rEI{3}}
  \closeproofpartindR

  \proofpartindR[Gewinnung einer zusammengesetzten Bijektion]{%
    A\EqCard B\dsep B\EqCard C
    \vdash \exists H\colon A\bij C%
  }{%
    A\EqCard B\dsep B\EqCard C
    \vdash \exists H\colon A\bij C%
  }[EqCardTransitiveBijectionExistence]
    \proofstep{1}{A\EqCard B}{\rA}
    \proofstep{2}{B\EqCard C}{\rA}
    \proofstep{1}{\exists F\colon A\bij B}{%
      \FormulaRefAuto{A\EqCard B\coloneqq
        \exists F\colon A\bij B}{1}}
    \proofstep{2}{\exists G\colon B\bij C}{%
      \FormulaRefAuto{A\EqCard B\coloneqq
        \exists F\colon A\bij B}{2}}
    \proofstep{5}{F\colon A\bij B}{\rA}
    \proofstep{6}{G\colon B\bij C}{\rA}
    \proofstep{5,6}{\exists H\colon A\bij C}{%
      \FormulaRefAuto{EqCardTransitiveFixedBijections}{5,6}}
    \proofstep{5}{\exists H\colon A\bij C}{\rEE{4,6,7}}
    \proofstep{1,2}{\exists H\colon A\bij C}{\rEE{3,5,8}}
  \closeproofpartindR

  \proofpartindR[Übergang zur Gleichmächtigkeit]{%
    A\EqCard B\dsep B\EqCard C\vdash A\EqCard C%
  }{%
    A\EqCard B\dsep B\EqCard C\vdash A\EqCard C%
  }[EqCardTransitiveConclusion]
    \proofstep{1}{A\EqCard B}{\rA}
    \proofstep{2}{B\EqCard C}{\rA}
    \proofstep{1,2}{\exists H\colon A\bij C}{%
      \FormulaRefAuto{EqCardTransitiveBijectionExistence}{1,2}}
    \proofstep{1,2}{A\EqCard C}{%
      \FormulaRefAuto{A\EqCard B\coloneqq
        \exists F\colon A\bij B}{3}}
  \closeproofpartindR
\end{tabproofsplit}

\subsection{Gleichmächtigkeit als Äquivalenzrelation}

Der Relationsoperator \(\EqCard\) ist bereits für beliebige Mengen definiert.
Auf jeder Mengenfamilie \(M\) wird er daher lediglich mit dem Träger
zusammen angegeben; eine zusätzliche Menge geordneter Paare ist nicht nötig.

\FormulaThmDeltaK[Gleichmächtigkeit ist eine Äquivalenzrelation]{%
  \EqRel{M,\EqCard}
}{EqCardRelEqRel}{%
  \DeltaRow{Mengen}{M\dsep A\dsep B\dsep C}
  \DeltaRow{Relationsoperatoren}{\EqCard}
}
\begin{tabproof}
  \proofstep{}{A\EqCard A}{\FormulaRefAuto{EqCardReflexive}}
  \proofstep{}{A\EqCard B\rightarrow B\EqCard A}{%
    \FormulaRefAuto{EqCardSymmetric}}
  \proofstep{}{%
    \bigl(A\EqCard B\land B\EqCard C\bigr)
    \rightarrow A\EqCard C}{\FormulaRefAuto{EqCardTransitive}}
  \proofstep{}{\EqRel{M,\EqCard}}{%
    \FormulaRefAuto{EqRelDef}{1,2,3}}
\end{tabproof}

\subsection{Weitere Folgerungen}

\FormulaThmDeltaK[Frische Adjunktion erhält Gleichmächtigkeit]{%
  A\EqCard B\dsep a\notin A\dsep b\notin B
  \vdash A\cup\{a\}\EqCard B\cup\{b\}%
}{EqCardFreshAdjunction}{%
  \DeltaRow{Mengen}{A\dsep B\dsep a\dsep b}%
  \DeltaRow{Funktionen}{F}%
}
\begin{tabproofsplit}
  \proofpartindR[Erweiterung einer festen Bijektion]{%
    \begin{aligned}[t]
      &a\notin A\dsep b\notin B\dsep F\colon A\bij B\vdash{}\\
      &\exists G\colon A\cup\{a\}\bij B\cup\{b\}
    \end{aligned}%
  }{%
    a\notin A\dsep b\notin B\dsep F\colon A\bij B
    \vdash \exists G\colon A\cup\{a\}\bij B\cup\{b\}%
  }[EqCardFreshAdjunctionFixedBijection]
    \proofstep{1}{a\notin A}{\rA}
    \proofstep{2}{b\notin B}{\rA}
    \proofstep{3}{F\colon A\bij B}{\rA}
    \proofstep{1,2,3}{%
      \ExtMap{\iota_{B,B\cup\{b\}}\circ F}{a}{b}
      \colon A\cup\{a\}\bij B\cup\{b\}}{%
      \FormulaRefAuto{FreshExtMapBijective}[thm]{3,1,2}}
    \proofstep{1,2,3}{%
      \exists G\colon A\cup\{a\}\bij B\cup\{b\}}{\rEI{4}}
  \closeproofpartindR

  \proofpartindR[Übergang zur Gleichmächtigkeit]{%
    A\EqCard B\dsep a\notin A\dsep b\notin B
    \vdash A\cup\{a\}\EqCard B\cup\{b\}%
  }{%
    A\EqCard B\dsep a\notin A\dsep b\notin B
    \vdash A\cup\{a\}\EqCard B\cup\{b\}%
  }[EqCardFreshAdjunctionConclusion]
    \proofstep{1}{A\EqCard B}{\rA}
    \proofstep{2}{a\notin A}{\rA}
    \proofstep{3}{b\notin B}{\rA}
    \proofstep{1}{\exists F\colon A\bij B}{%
      \FormulaRefAuto{A\EqCard B\coloneqq
        \exists F\colon A\bij B}{1}}
    \proofstep{5}{F\colon A\bij B}{\rA}
    \proofstep{2,3,5}{%
      \exists G\colon A\cup\{a\}\bij B\cup\{b\}}{%
      \FormulaRefAuto{EqCardFreshAdjunctionFixedBijection}{2,3,5}}
    \proofstep{2,3,5}{A\cup\{a\}\EqCard B\cup\{b\}}{%
      \FormulaRefAuto{A\EqCard B\coloneqq
        \exists F\colon A\bij B}{6}}
    \proofstep{1,2,3}{A\cup\{a\}\EqCard B\cup\{b\}}{%
      \rEE{4,5,7}}
  \closeproofpartindR
\end{tabproofsplit}

\FormulaThmDeltaK[Transport von Injektionen entlang der Gleichmächtigkeit]{%
  A \EqCard M\dsep B \EqCard N\dsep F\colon A\inj B
  \vdash \exists G\colon M\inj N
}{EqCardInjectionTransport}{%
  \DeltaRow{Mengen}{A\dsep B\dsep M\dsep N}
  \DeltaRow{Funktionen}{F\dsep G\dsep H\dsep K}
}
\begin{tabproofsplitwide}
  \proofpartwideindR[Transport entlang fester Bijektionen]{%
    F\colon A\inj B\dsep H\colon A\bij M\dsep K\colon B\bij N
    \vdash \exists G\colon M\inj N
  }{%
    F\colon A\inj B\dsep H\colon A\bij M\dsep K\colon B\bij N
    \vdash \exists G\colon M\inj N
  }[EqCardInjectionTransportFixedBijections]
    \proofstepwidestar[1]{F\colon A\inj B}{\rA}
    \proofstepwidestar[2]{H\colon A\bij M}{\rA}
    \proofstepwidestar[3]{K\colon B\bij N}{\rA}
    \proofstepwidestar[2]{H^{-1}\colon M\bij A}{%
      \FormulaRefAuto{F\colon A\bij B\vdash F^{-1}\colon B\bij A}{2}}
    \proofstepwidestar[2]{H^{-1}\colon M\inj A}{%
      \FormulaRefAuto{F\colon A\bij B \vdash F\colon A\inj B}{4}}
    \proofstepwidestar[1,2]{F\circ H^{-1}\colon M\inj B}{%
      \FormulaRefAuto{F\colon A\inj B \dsep G\colon B\inj C
      \vdash G\circ F\colon A\inj C}{5,1}}
    \proofstepwidestar[3]{K\colon B\inj N}{%
      \FormulaRefAuto{F\colon A\bij B \vdash F\colon A\inj B}{3}}
    \proofstepwidestar[1,2,3]{K\circ(F\circ H^{-1})\colon M\inj N}{%
      \FormulaRefAuto{F\colon A\inj B \dsep G\colon B\inj C
      \vdash G\circ F\colon A\inj C}{6,7}}
    \proofstepwidestar[1,2,3]{\exists G\colon M\inj N}{\rEI{8}}
  \closeproofpartwideindR

  \proofpartwide{Schluss}
    \proofstepwidestar[1]{A \EqCard M}{\rA}
    \proofstepwidestar[2]{B \EqCard N}{\rA}
    \proofstepwidestar[3]{F\colon A\inj B}{\rA}
    \proofstepwidestar[1]{\exists H\colon A\bij M}{%
      \FormulaRefAuto{A \EqCard B \coloneqq \exists F\colon A\bij B}{1}}
    \proofstepwidestar[2]{\exists K\colon B\bij N}{%
      \FormulaRefAuto{A \EqCard B \coloneqq \exists F\colon A\bij B}{2}}
    \proofstepwidestar[6]{H\colon A\bij M}{\rA}
    \proofstepwidestar[7]{K\colon B\bij N}{\rA}
    \proofstepwidestar[3,6,7]{\exists G\colon M\inj N}{%
      \FormulaRefAuto{EqCardInjectionTransportFixedBijections}{3,6,7}}
    \proofstepwidestar[2,3,6]{\exists G\colon M\inj N}{%
      \rEE{5,7,8}}
    \proofstepwidestar[1,2,3]{\exists G\colon M\inj N}{%
      \rEE{4,6,9}}
  \closeproofpartwide
\end{tabproofsplitwide}


\FormulaThmDeltaK[Potenzmengen gleichmächtiger Mengen sind gleichmächtig]{%
  A \EqCard B \vdash \powerset(A)\EqCard \powerset(B)
}{EqCardPowerSetForward}{%
  \DeltaRow{Mengen}{A\dsep B}
  \DeltaRow{Funktionen}{F\dsep G\dsep \widehat{F}}
}
\begin{tabproof}
  \proofstep{1}{A \EqCard B}{\rA}
  \proofstep{1}{\exists F\colon A\bij B}{%
    \FormulaRefAuto{A \EqCard B \coloneqq \exists F\colon A\bij B}{1}}

  \proofstep{3}{F\colon A\bij B}{\rA}
    \proofstep{3}{\PowMap{F}\colon \powerset(A)\bij \powerset(B)}{%
      \FormulaRefAuto{F\colon A\bij B \vdash \PowMap{F}\colon \powerset(A)\bij \powerset(B)}{3}}
  \proofstep{3}{\exists G\colon \powerset(A)\bij \powerset(B)}{\rEI{4}}

  \proofstep{1}{\exists G\colon \powerset(A)\bij \powerset(B)}{\rEE{2,3,5}}
  \proofstep{1}{\powerset(A)\EqCard \powerset(B)}{%
    \FormulaRefAuto{A \EqCard B \coloneqq \exists F\colon A\bij B}{6}}
\end{tabproof}

\FormulaThmDeltaK[Gleichmächtigkeit und ein induzierter
Potenzmengenzeuge]{%
  \begin{aligned}[t]
    A\EqCard B\quad\leftrightarrow\quad
    \exists F\,\bigl(&F\colon A\to B\land{}\\[-2pt]
      &\PowMap{F}\colon\powerset(A)\bij\powerset(B)\bigr)
  \end{aligned}%
}{EqCardPowerMapWitnessEquiv}{%
  \DeltaRow{Mengen}{A\dsep B}
  \DeltaRow{Funktionen}{F\dsep\PowMap{F}}
}
\begin{tabproofsplitwide}
  \proofpartwideindR[Hinrichtung]{%
    A\EqCard B\vdash
    \exists F\,\bigl(F\colon A\to B\land
      \PowMap{F}\colon\powerset(A)\bij\powerset(B)\bigr)%
  }{%
    A\EqCard B\vdash
    \exists F\,\bigl(F\colon A\to B\land
      \PowMap{F}\colon\powerset(A)\bij\powerset(B)\bigr)%
  }[EqCardPowerMapWitnessForward]
    \proofstepwidestar[1]{A\EqCard B}{\rA}
    \proofstepwidestar[1]{\exists F\colon A\bij B}{%
      \FormulaRefAuto{A\EqCard B\coloneqq\exists F\colon A\bij B}{1}}
    \proofstepwidestar[3]{F\colon A\bij B}{\rA}
    \proofstepwidestar[3]{F\colon A\to B}{%
      \FormulaRefAuto{Bijektive Funktion}[def]{3}}
    \proofstepwidestar[3]{%
      \PowMap{F}\colon\powerset(A)\bij\powerset(B)}{%
      \FormulaRefAuto{F\colon A\bij B\vdash
        \PowMap{F}\colon\powerset(A)\bij\powerset(B)}[thm]{3}}
    \proofstepwidestar[3]{%
      F\colon A\to B\land
      \PowMap{F}\colon\powerset(A)\bij\powerset(B)}{\rAI{4,5}}
    \proofstepwidestar[3]{%
      \exists F\,\bigl(F\colon A\to B\land
        \PowMap{F}\colon\powerset(A)\bij\powerset(B)\bigr)}{\rEI{6}}
    \proofstepwidestar[1]{%
      \exists F\,\bigl(F\colon A\to B\land
        \PowMap{F}\colon\powerset(A)\bij\powerset(B)\bigr)}{%
      \rEE{2,3,7}}
  \closeproofpartwideindR

  \proofpartwideindR[Rückrichtung]{%
    \begin{aligned}[t]
      &\exists F\,\bigl(F\colon A\to B\land
        \PowMap{F}\colon\powerset(A)\bij\powerset(B)\bigr)\vdash{}\\[-2pt]
      &A\EqCard B
    \end{aligned}%
  }{%
    \exists F\,\bigl(F\colon A\to B\land
      \PowMap{F}\colon\powerset(A)\bij\powerset(B)\bigr)
    \vdash A\EqCard B%
  }[EqCardPowerMapWitnessBackward]
    \proofstepwidestar[1]{%
      \exists F\,\bigl(F\colon A\to B\land
        \PowMap{F}\colon\powerset(A)\bij\powerset(B)\bigr)}{\rA}
    \proofstepwidestar[2]{%
      F\colon A\to B\land
      \PowMap{F}\colon\powerset(A)\bij\powerset(B)}{\rA}
    \proofstepwidestar[2]{F\colon A\to B}{\rAEa{2}}
    \proofstepwidestar[2]{%
      \PowMap{F}\colon\powerset(A)\bij\powerset(B)}{\rAEb{2}}
    \proofstepwidestar[2]{F\colon A\bij B}{%
      \FormulaRefAuto{PowerMapBijectiveReflectsBijectivity}[thm]{3,4}}
    \proofstepwidestar[2]{\exists G\colon A\bij B}{\rEI{5}}
    \proofstepwidestar[2]{A\EqCard B}{%
      \FormulaRefAuto{A\EqCard B\coloneqq\exists F\colon A\bij B}{6}}
    \proofstepwidestar[1]{A\EqCard B}{\rEE{1,2,7}}
  \closeproofpartwideindR

  \proofpartwide{Zusammenführung}
    \proofstepwidestar[]{%
      A\EqCard B\rightarrow
      \exists F\,\bigl(F\colon A\to B\land
        \PowMap{F}\colon\powerset(A)\bij\powerset(B)\bigr)}{%
      \FormulaRefAuto{EqCardPowerMapWitnessForward}}
    \proofstepwidestar[]{%
      \exists F\,\bigl(F\colon A\to B\land
        \PowMap{F}\colon\powerset(A)\bij\powerset(B)\bigr)
      \rightarrow A\EqCard B}{%
      \FormulaRefAuto{EqCardPowerMapWitnessBackward}}
    \proofstepwidestar[]{%
      \begin{aligned}[t]
        A\EqCard B\quad\leftrightarrow\quad
        \exists F\,\bigl(&F\colon A\to B\land{}\\[-2pt]
          &\PowMap{F}\colon\powerset(A)\bij\powerset(B)\bigr)
      \end{aligned}}{\rLRI{1,2}}
  \closeproofpartwide
\end{tabproofsplitwide}

\FormulaThmDeltaK[Induzierte Restfamilienbijektion erzwingt Gleichmächtigkeit]{%
  \begin{aligned}[t]
    &\forall a\in A\,\bigl(\{a\}\notin C\bigr)\dsep
      \forall b\in B\,\bigl(\{b\}\notin C\bigr)\dsep{}\\[-2pt]
    &\exists F\,\Bigl(F\colon A\to B\land
      \PowMap{F}\restriction_{\powerset(A)\setminus C}
      \colon\powerset(A)\setminus C\bij\powerset(B)\setminus C\Bigr)
      \vdash A\EqCard B
  \end{aligned}%
}{EqCardRestrictedPowerMapWitness}{%
  \DeltaRow{Mengen}{A\dsep B\dsep C}
  \DeltaRow{Funktionen}{F\dsep\PowMap{F}}
}
\begin{tabproofwide}
  \proofstepwidestar[1]{\forall a\in A\,(\{a\}\notin C)}{\rA}
  \proofstepwidestar[2]{\forall b\in B\,(\{b\}\notin C)}{\rA}
  \proofstepwidestarbreakkeep[3]{%
    \exists F\,\Bigl(F\colon A\to B\land
      \PowMap{F}\restriction_{\powerset(A)\setminus C}
      \colon\powerset(A)\setminus C\bij\powerset(B)\setminus C\Bigr)}{\rA}
  \proofstepwidestarbreakkeep[4]{%
    F\colon A\to B\land
    \PowMap{F}\restriction_{\powerset(A)\setminus C}
      \colon\powerset(A)\setminus C\bij\powerset(B)\setminus C}{\rA}
  \proofstepwidestar[4]{F\colon A\to B}{\rAEa{4}}
  \proofstepwidestarbreakkeep[4]{%
    \PowMap{F}\restriction_{\powerset(A)\setminus C}
      \colon\powerset(A)\setminus C\bij\powerset(B)\setminus C}{\rAEb{4}}
  \proofstepwidestar[1,2,4]{F\colon A\bij B}{%
    \FormulaRefAuto{RestrictedPowerMapBijectiveReflectsBijectivity}[thm]{%
      5,1,2,6}}
  \proofstepwidestar[1,2,4]{\exists G\colon A\bij B}{\rEI{7}}
  \proofstepwidestar[1,2,4]{A\EqCard B}{%
    \FormulaRefAuto{A\EqCard B\coloneqq\exists F\colon A\bij B}{8}}
  \proofstepwidestar[1,2,3]{A\EqCard B}{\rEE{3,4,9}}
\end{tabproofwide}

\begin{remark}[Gleichmächtige Potenzmengen lassen sich nicht kürzen]
Die Aussage
\[
  \powerset(A)\EqCard\powerset(B)\ \longrightarrow\ A\EqCard B
\]
ist nicht die Rückrichtung des vorstehenden Satzes: Auf der linken Seite
darf eine beliebige Potenzmengenbijektion als Zeuge auftreten, während dort
ausdrücklich ein von \(F\) induzierter Zeuge verlangt wird.  Die behauptete
Kürzungsregel ist, unter der üblichen relativen Konsistenzvoraussetzung, von
ZFC unabhängig: Es gibt Modelle dieser Mengenlehre, in denen verschieden
mächtige unendliche Mengen gleichmächtige Potenzmengen besitzen, und Modelle,
in denen dies nicht vorkommt.  Der vorangehende Band kennzeichnet diesen
externen mengentheoretischen Befund ausdrücklich und grenzt ihn von der
bewiesenen induzierten Rückwirkung ab.

Auch
\[
  \powerset(A)\setminus C\EqCard\powerset(B)\setminus C
  \ \longrightarrow\ A\EqCard B
\]
ist für beliebiges \(C\) falsch.  Für
\(A=\varnothing\), \(B=\{\varnothing\}\) und
\(C=\{B\}=\{\{\varnothing\}\}\) sind beide Restfamilien gleich
\(\{\varnothing\}\), obwohl die leere und die nichtleere Grundmenge nicht
gleichmächtig sind.  Der Satz
\FormulaRefAuto{EqCardRestrictedPowerMapWitness}[thm] zeigt die korrekte
induzierte Fassung mit hinreichenden Einermengenbedingungen.
\end{remark}

\chapter{Äquivalenzklassen und Quotientenmengen}

\section{Äquivalenzklassen}

\FormulaDefDeltaK[Äquivalenzklasse eines Elements]{%
  \EqRel{A,\sim}\dsep x\in A\vdash
  \EqClass{x}{\sim}{A}\coloneqq
  \{y\in A\mid x\sim y\}
}{EqClassDef}{%
  \DeltaRow{Mengen}{A\dsep x\dsep y}
  \DeltaRow{Relationsoperatoren}{\sim}
  \DeltaRow{Neue Symbole}{\EqClass{x}{\sim}{A}}
}

\FormulaThmDeltaK[Charakterisierung der Äquivalenzklasse]{%
  \EqRel{A,\sim}\dsep x\in A
  \vdash
  y\in \EqClass{x}{\sim}{A}
  \leftrightarrow
  y\in A\land x\sim y
}{EqClassChar}{%
  \DeltaRow{Mengen}{A\dsep x\dsep y}
  \DeltaRow{Relationsoperatoren}{\sim}
}
\begin{tabproofsplit}
  \proofpartindR[Hinrichtung]{%
    \EqRel{A,\sim}\dsep x\in A\vdash
    y\in \EqClass{x}{\sim}{A}
    \rightarrow y\in A\land x\sim y%
  }{%
    \EqRel{A,\sim}\dsep x\in A\vdash
    y\in \EqClass{x}{\sim}{A}
    \rightarrow y\in A\land x\sim y%
  }[EqClassCharForwardDirection]
    \proofstep{1}{\EqRel{A,\sim}}{\rA}
    \proofstep{2}{x\in A}{\rA}
    \proofstep{1,2}{%
      \EqClass{x}{\sim}{A}\coloneqq
      \{y\in A\mid x\sim y\}}{%
      \FormulaRefAuto{EqClassDef}{1,2}}
    \proofstep{4}{y\in \EqClass{x}{\sim}{A}}{\rA}
    \proofstep{1,2,4}{y\in A\land x\sim y}{%
      \FormulaRefAuto{x \in \{u \in A \mid P(u)\} \eqvdash
        x \in A \land P(x)}{3,4}}
    \proofstep{1,2}{%
      y\in \EqClass{x}{\sim}{A}
      \rightarrow y\in A\land x\sim y}{\rRI{4,5}}
  \closeproofpartindR

  \proofpartindR[Rückrichtung]{%
    \EqRel{A,\sim}\dsep x\in A\vdash
    y\in A\land x\sim y
    \rightarrow y\in \EqClass{x}{\sim}{A}%
  }{%
    \EqRel{A,\sim}\dsep x\in A\vdash
    y\in A\land x\sim y
    \rightarrow y\in \EqClass{x}{\sim}{A}%
  }[EqClassCharBackwardDirection]
    \proofstep{1}{\EqRel{A,\sim}}{\rA}
    \proofstep{2}{x\in A}{\rA}
    \proofstep{1,2}{%
      \EqClass{x}{\sim}{A}\coloneqq
      \{y\in A\mid x\sim y\}}{%
      \FormulaRefAuto{EqClassDef}{1,2}}
    \proofstep{4}{y\in A\land x\sim y}{\rA}
    \proofstep{1,2,4}{y\in \EqClass{x}{\sim}{A}}{%
      \FormulaRefAuto{x \in \{u \in A \mid P(u)\} \eqvdash
        x \in A \land P(x)}{3,4}}
    \proofstep{1,2}{%
      y\in A\land x\sim y
      \rightarrow y\in \EqClass{x}{\sim}{A}}{\rRI{4,5}}
  \closeproofpartindR

  \proofpartindR[Zusammenführung]{%
    \EqRel{A,\sim}\dsep x\in A\vdash
    y\in \EqClass{x}{\sim}{A}
    \leftrightarrow y\in A\land x\sim y%
  }{%
    \EqRel{A,\sim}\dsep x\in A\vdash
    y\in \EqClass{x}{\sim}{A}
    \leftrightarrow y\in A\land x\sim y%
  }[EqClassCharEquivalenceConclusion]
    \proofstep{1}{\EqRel{A,\sim}}{\rA}
    \proofstep{2}{x\in A}{\rA}
    \proofstep{1,2}{%
      y\in \EqClass{x}{\sim}{A}
      \rightarrow y\in A\land x\sim y}{%
      \FormulaRefAuto{EqClassCharForwardDirection}{1,2}}
    \proofstep{1,2}{%
      y\in A\land x\sim y
      \rightarrow y\in \EqClass{x}{\sim}{A}}{%
      \FormulaRefAuto{EqClassCharBackwardDirection}{1,2}}
    \proofstep{1,2}{%
      y\in \EqClass{x}{\sim}{A}
      \leftrightarrow y\in A\land x\sim y}{\rLRI{3,4}}
  \closeproofpartindR
\end{tabproofsplit}

\FormulaThmDeltaK[Repräsentant liegt in seiner Äquivalenzklasse]{%
  \EqRel{A,\sim}\dsep x\in A
  \vdash
  x\in \EqClass{x}{\sim}{A}
}{EqClassRepresentative}{%
  \DeltaRow{Mengen}{A\dsep x}
  \DeltaRow{Relationsoperatoren}{\sim}
}
\begin{tabproof}
  \proofstep{1}{\EqRel{A,\sim}}{\rA}
  \proofstep{2}{x\in A}{\rA}
  \proofstep{1,2}{x\sim x}{\FormulaRefAuto{EqRelDef}{1,2}}
  \proofstep{1,2}{x\in A\land x\sim x}{\rAI{2,3}}
  \proofstep{1,2}{x\in \EqClass{x}{\sim}{A}\leftrightarrow
    x\in A\land x\sim x}{\FormulaRefAuto{EqClassChar}{1,2}}
  \proofstep{1,2}{x\in \EqClass{x}{\sim}{A}}{%
    \FormulaRefAuto{P \leftrightarrow Q, Q \vdash P}{5,4}}
\end{tabproof}

\FormulaThmDeltaK[Elemente einer Äquivalenzklasse liegen im Grundbereich]{%
  \EqRel{A,\sim}\dsep x\in A\dsep y\in \EqClass{x}{\sim}{A}
  \vdash y\in A
}{EqClassElementInCarrier}{%
  \DeltaRow{Mengen}{A\dsep x\dsep y}
  \DeltaRow{Relationsoperatoren}{\sim}
}
\begin{tabproof}
  \proofstep{1}{\EqRel{A,\sim}}{\rA}
  \proofstep{2}{x\in A}{\rA}
  \proofstep{3}{y\in \EqClass{x}{\sim}{A}}{\rA}
  \proofstep{1,2}{y\in \EqClass{x}{\sim}{A}\leftrightarrow
    y\in A\land x\sim y}{\FormulaRefAuto{EqClassChar}{1,2}}
  \proofstep{1,2,3}{y\in A\land x\sim y}{%
    \FormulaRefAuto{P \leftrightarrow Q, P \vdash Q}{4,3}}
  \proofstep{1,2,3}{y\in A}{\rAEa{5}}
\end{tabproof}

\FormulaThmDeltaK[Elemente einer Äquivalenzklasse sind äquivalent zum Repräsentanten]{%
  \EqRel{A,\sim}\dsep x\in A\dsep y\in \EqClass{x}{\sim}{A}
  \vdash x\sim y
}{EqClassElementRelated}{%
  \DeltaRow{Mengen}{A\dsep x\dsep y}
  \DeltaRow{Relationsoperatoren}{\sim}
}
\begin{tabproof}
  \proofstep{1}{\EqRel{A,\sim}}{\rA}
  \proofstep{2}{x\in A}{\rA}
  \proofstep{3}{y\in \EqClass{x}{\sim}{A}}{\rA}
  \proofstep{1,2}{y\in \EqClass{x}{\sim}{A}\leftrightarrow
    y\in A\land x\sim y}{\FormulaRefAuto{EqClassChar}{1,2}}
  \proofstep{1,2,3}{y\in A\land x\sim y}{%
    \FormulaRefAuto{P \leftrightarrow Q, P \vdash Q}{4,3}}
  \proofstep{1,2,3}{x\sim y}{\rAEb{5}}
\end{tabproof}

\FormulaThmDeltaK[Äquivalente Elemente haben dieselbe Äquivalenzklasse]{%
  \EqRel{A,\sim}\dsep x\in A\dsep y\in A\dsep x\sim y
  \vdash
  \EqClass{x}{\sim}{A}=\EqClass{y}{\sim}{A}
}{EqClassEqualFromRel}{%
  \DeltaRow{Mengen}{A\dsep x\dsep y\dsep z}
  \DeltaRow{Relationsoperatoren}{\sim}
}
\begin{tabproofsplitwide}
  \proofpartwideindR[Teilmengenrichtung \(\EqClass{x}{\sim}{A}\subseteq\EqClass{y}{\sim}{A}\)]{%
    \EqRel{A,\sim}\dsep x\in A\dsep y\in A\dsep x\sim y
    \vdash
    \EqClass{x}{\sim}{A}\subseteq\EqClass{y}{\sim}{A}
  }{%
    \EqRel{A,\sim}\dsep x\in A\dsep y\in A\dsep x\sim y
    \vdash
    \EqClass{x}{\sim}{A}\subseteq\EqClass{y}{\sim}{A}
  }[EqClassEqualFromRelSubsetLeft]
    \proofstepwidestar[1]{\EqRel{A,\sim}}{\rA}
    \proofstepwidestar[2]{x\in A}{\rA}
    \proofstepwidestar[3]{y\in A}{\rA}
    \proofstepwidestar[4]{x\sim y}{\rA}
    \proofstepwidestar[5]{z\in \EqClass{x}{\sim}{A}}{\rA}
    \proofstepwidestar[1,2,5]{z\in A}{%
      \FormulaRefAuto{EqClassElementInCarrier}{1,2,5}}
    \proofstepwidestar[1,2,5]{x\sim z}{%
      \FormulaRefAuto{EqClassElementRelated}{1,2,5}}
    \proofstepwidestar[1,4]{y\sim x}{\FormulaRefAuto{EqRelSymmetry}{4}}
    \proofstepwidestar[1,2,4,5]{y\sim z}{%
      \FormulaRefAuto{x\sim y \dsep y\sim z \vdash x\sim z}[ax]{8,7}}
    \proofstepwidestar[1,2,3,4,5]{z\in A\land y\sim z}{\rAI{6,9}}
    \proofstepwidestar[1,3]{z\in \EqClass{y}{\sim}{A}\leftrightarrow
      z\in A\land y\sim z}{\FormulaRefAuto{EqClassChar}{1,3}}
    \proofstepwidestar[1,2,3,4,5]{z\in \EqClass{y}{\sim}{A}}{%
      \FormulaRefAuto{P \leftrightarrow Q, Q \vdash P}{11,10}}
    \proofstepwidestar[1,2,3,4]{\EqClass{x}{\sim}{A}\subseteq\EqClass{y}{\sim}{A}}{%
      \FormulaRefAuto{A\subseteq B \coloneq
      \forall x\,(x\in A \rightarrow x\in B)}{\rUI{\rRI{5,12}}}}
  \closeproofpartwideind

  \proofpartwideindR[Teilmengenrichtung \(\EqClass{y}{\sim}{A}\subseteq\EqClass{x}{\sim}{A}\)]{%
    \EqRel{A,\sim}\dsep x\in A\dsep y\in A\dsep x\sim y
    \vdash
    \EqClass{y}{\sim}{A}\subseteq\EqClass{x}{\sim}{A}
  }{%
    \EqRel{A,\sim}\dsep x\in A\dsep y\in A\dsep x\sim y
    \vdash
    \EqClass{y}{\sim}{A}\subseteq\EqClass{x}{\sim}{A}
  }[EqClassEqualFromRelSubsetRight]
    \proofstepwidestar[1]{\EqRel{A,\sim}}{\rA}
    \proofstepwidestar[2]{x\in A}{\rA}
    \proofstepwidestar[3]{y\in A}{\rA}
    \proofstepwidestar[4]{x\sim y}{\rA}
    \proofstepwidestar[5]{z\in \EqClass{y}{\sim}{A}}{\rA}
    \proofstepwidestar[1,3,5]{z\in A}{%
      \FormulaRefAuto{EqClassElementInCarrier}{1,3,5}}
    \proofstepwidestar[1,3,5]{y\sim z}{%
      \FormulaRefAuto{EqClassElementRelated}{1,3,5}}
    \proofstepwidestar[1,3,4,5]{x\sim z}{%
      \FormulaRefAuto{x\sim y \dsep y\sim z \vdash x\sim z}[ax]{4,7}}
    \proofstepwidestar[1,2,3,4,5]{z\in A\land x\sim z}{\rAI{6,8}}
    \proofstepwidestar[1,2]{z\in \EqClass{x}{\sim}{A}\leftrightarrow
      z\in A\land x\sim z}{\FormulaRefAuto{EqClassChar}{1,2}}
    \proofstepwidestar[1,2,3,4,5]{z\in \EqClass{x}{\sim}{A}}{%
      \FormulaRefAuto{P \leftrightarrow Q, Q \vdash P}{10,9}}
    \proofstepwidestar[1,2,3,4]{\EqClass{y}{\sim}{A}\subseteq\EqClass{x}{\sim}{A}}{%
      \FormulaRefAuto{A\subseteq B \coloneq
      \forall x\,(x\in A \rightarrow x\in B)}{\rUI{\rRI{5,11}}}}
  \closeproofpartwideind

  \proofpartwide{Schluss}
    \proofstepwidestar[1]{\EqRel{A,\sim}}{\rA}
    \proofstepwidestar[2]{x\in A}{\rA}
    \proofstepwidestar[3]{y\in A}{\rA}
    \proofstepwidestar[4]{x\sim y}{\rA}
    \proofstepwidestar[1,2,3,4]{\EqClass{x}{\sim}{A}\subseteq\EqClass{y}{\sim}{A}}{%
      \FormulaRefAuto{EqClassEqualFromRelSubsetLeft}{1,2,3,4}}
    \proofstepwidestar[1,2,3,4]{\EqClass{y}{\sim}{A}\subseteq\EqClass{x}{\sim}{A}}{%
      \FormulaRefAuto{EqClassEqualFromRelSubsetRight}{1,2,3,4}}
    \proofstepwidestar[1,2,3,4]{\EqClass{x}{\sim}{A}=\EqClass{y}{\sim}{A}}{%
      \FormulaRefAuto{A \subseteq B, B \subseteq A \vdash A = B}{5,6}}
  \closeproofpartwide
\end{tabproofsplitwide}

\FormulaThmDeltaK[Gleiche Äquivalenzklassen liefern Äquivalenz]{%
  \EqRel{A,\sim}\dsep x\in A\dsep y\in A\dsep
  \EqClass{x}{\sim}{A}=\EqClass{y}{\sim}{A}
  \vdash
  x\sim y
}{EqClassRelFromEqual}{%
  \DeltaRow{Mengen}{A\dsep x\dsep y}
  \DeltaRow{Relationsoperatoren}{\sim}
}
\begin{tabproof}
  \proofstep{1}{\EqRel{A,\sim}}{\rA}
  \proofstep{2}{x\in A}{\rA}
  \proofstep{3}{y\in A}{\rA}
  \proofstep{4}{\EqClass{x}{\sim}{A}=\EqClass{y}{\sim}{A}}{\rA}
  \proofstep{1,2}{x\in \EqClass{x}{\sim}{A}}{%
    \FormulaRefAuto{EqClassRepresentative}{1,2}}
  \proofstep{1,2,4}{x\in \EqClass{y}{\sim}{A}}{\rIE{4,5}}
  \proofstep{1,2,3,4}{y\sim x}{%
    \FormulaRefAuto{EqClassElementRelated}{1,3,6}}
  \proofstep{1,2,3,4}{x\sim y}{\FormulaRefAuto{EqRelSymmetry}{7}}
\end{tabproof}

\FormulaThmDeltaK[Charakterisierung gleicher Äquivalenzklassen]{%
  \EqRel{A,\sim}\dsep x\in A\dsep y\in A
  \vdash
  \bigl(\EqClass{x}{\sim}{A}=\EqClass{y}{\sim}{A}\bigr)
  \leftrightarrow
  x\sim y
}{EqClassEqualIffRel}{%
  \DeltaRow{Mengen}{A\dsep x\dsep y}
  \DeltaRow{Relationsoperatoren}{\sim}
}
\begin{tabproofsplit}
  \proofpartindR[Von gleichen Klassen zur Relation]{%
    \begin{aligned}[t]
      &\EqRel{A,\sim}\dsep x\in A\dsep y\in A\vdash{}\\
      &\EqClass{x}{\sim}{A}=\EqClass{y}{\sim}{A}
      \rightarrow x\sim y
    \end{aligned}%
  }{%
    \EqRel{A,\sim}\dsep x\in A\dsep y\in A\vdash
    \EqClass{x}{\sim}{A}=\EqClass{y}{\sim}{A}
    \rightarrow x\sim y%
  }[EqClassEqualIffRelForwardDirection]
    \proofstep{1}{\EqRel{A,\sim}}{\rA}
    \proofstep{2}{x\in A}{\rA}
    \proofstep{3}{y\in A}{\rA}
    \proofstep{4}{\EqClass{x}{\sim}{A}=\EqClass{y}{\sim}{A}}{\rA}
    \proofstep{1,2,3,4}{x\sim y}{%
      \FormulaRefAuto{EqClassRelFromEqual}{1,2,3,4}}
    \proofstep{1,2,3}{%
      \EqClass{x}{\sim}{A}=\EqClass{y}{\sim}{A}
      \rightarrow x\sim y}{\rRI{4,5}}
  \closeproofpartindR

  \proofpartindR[Von der Relation zu gleichen Klassen]{%
    \begin{aligned}[t]
      &\EqRel{A,\sim}\dsep x\in A\dsep y\in A\vdash{}\\
      &x\sim y
      \rightarrow \EqClass{x}{\sim}{A}=\EqClass{y}{\sim}{A}
    \end{aligned}%
  }{%
    \EqRel{A,\sim}\dsep x\in A\dsep y\in A\vdash
    x\sim y
    \rightarrow \EqClass{x}{\sim}{A}=\EqClass{y}{\sim}{A}%
  }[EqClassEqualIffRelBackwardDirection]
    \proofstep{1}{\EqRel{A,\sim}}{\rA}
    \proofstep{2}{x\in A}{\rA}
    \proofstep{3}{y\in A}{\rA}
    \proofstep{4}{x\sim y}{\rA}
    \proofstep{1,2,3,4}{\EqClass{x}{\sim}{A}=\EqClass{y}{\sim}{A}}{%
      \FormulaRefAuto{EqClassEqualFromRel}{1,2,3,4}}
    \proofstep{1,2,3}{%
      x\sim y
      \rightarrow \EqClass{x}{\sim}{A}=\EqClass{y}{\sim}{A}}{\rRI{4,5}}
  \closeproofpartindR

  \proofpartindR[Zusammenführung]{%
    \begin{aligned}[t]
      &\EqRel{A,\sim}\dsep x\in A\dsep y\in A\vdash{}\\
      &\bigl(\EqClass{x}{\sim}{A}=\EqClass{y}{\sim}{A}\bigr)
      \leftrightarrow x\sim y
    \end{aligned}%
  }{%
    \EqRel{A,\sim}\dsep x\in A\dsep y\in A\vdash
    \bigl(\EqClass{x}{\sim}{A}=\EqClass{y}{\sim}{A}\bigr)
    \leftrightarrow x\sim y%
  }[EqClassEqualIffRelConclusion]
    \proofstep{1}{\EqRel{A,\sim}}{\rA}
    \proofstep{2}{x\in A}{\rA}
    \proofstep{3}{y\in A}{\rA}
    \proofstep{1,2,3}{%
      \EqClass{x}{\sim}{A}=\EqClass{y}{\sim}{A}
      \rightarrow x\sim y}{%
      \FormulaRefAuto{EqClassEqualIffRelForwardDirection}{1,2,3}}
    \proofstep{1,2,3}{%
      x\sim y
      \rightarrow \EqClass{x}{\sim}{A}=\EqClass{y}{\sim}{A}}{%
      \FormulaRefAuto{EqClassEqualIffRelBackwardDirection}{1,2,3}}
    \proofstep{1,2,3}{%
      \bigl(\EqClass{x}{\sim}{A}=\EqClass{y}{\sim}{A}\bigr)
      \leftrightarrow x\sim y}{\rLRI{4,5}}
  \closeproofpartindR
\end{tabproofsplit}

\section{Quotientenmengen}

\FormulaDefDeltaK[Quotientenmenge einer Äquivalenzrelation]{%
  \EqRel{A,\sim}\vdash
  \QuotientSet{A}{\sim}\coloneqq
  \{C\in\powerset(A)\mid
  \exists x\in A\,C=\EqClass{x}{\sim}{A}\}
}{QuotientSetDef}{%
  \DeltaRow{Mengen}{A\dsep C\dsep x}
  \DeltaRow{Relationsoperatoren}{\sim}
  \DeltaRow{Äquivalenzklassen}{\EqClass{x}{\sim}{A}}
  \DeltaRow{Neue Symbole}{\QuotientSet{A}{\sim}}
}

\FormulaThmDeltaK[Charakterisierung der Quotientenmenge]{%
  \EqRel{A,\sim}
  \vdash
  C\in \QuotientSet{A}{\sim}
  \leftrightarrow
  C\in\powerset(A)\land
  \exists x\in A\,C=\EqClass{x}{\sim}{A}
}{QuotientSetChar}{%
  \DeltaRow{Mengen}{A\dsep C\dsep x}
  \DeltaRow{Relationsoperatoren}{\sim}
  \DeltaRow{Äquivalenzklassen}{\EqClass{x}{\sim}{A}}
}
\begin{tabproof}
  \proofstep{1}{\EqRel{A,\sim}}{\rA}
  \proofstep{1}{%
    \QuotientSet{A}{\sim}\coloneqq
    \{C\in\powerset(A)\mid
    \exists x\in A\,C=\EqClass{x}{\sim}{A}\}}{%
    \FormulaRefAuto{QuotientSetDef}{1}}

  \proofstep{3}{C\in \QuotientSet{A}{\sim}}{\rA}
  \proofstep{1,3}{C\in\powerset(A)\land
    \exists x\in A\,C=\EqClass{x}{\sim}{A}}{%
    \FormulaRefAuto{x \in \{u \in A \mid P(u)\} \eqvdash x \in A \land P(x)}{2,3}}
  \proofstep{1}{C\in \QuotientSet{A}{\sim}
    \rightarrow C\in\powerset(A)\land
    \exists x\in A\,C=\EqClass{x}{\sim}{A}}{\rRI{3,4}}

  \proofstep{6}{C\in\powerset(A)\land
    \exists x\in A\,C=\EqClass{x}{\sim}{A}}{\rA}
  \proofstep{1,6}{C\in \QuotientSet{A}{\sim}}{%
    \FormulaRefAuto{x \in \{u \in A \mid P(u)\} \eqvdash x \in A \land P(x)}{2,6}}
  \proofstep{1}{C\in\powerset(A)\land
    \exists x\in A\,C=\EqClass{x}{\sim}{A}
    \rightarrow C\in \QuotientSet{A}{\sim}}{\rRI{6,7}}

  \proofstep{1}{C\in \QuotientSet{A}{\sim}
    \leftrightarrow
    C\in\powerset(A)\land
    \exists x\in A\,C=\EqClass{x}{\sim}{A}}{\rLRI{5,8}}
\end{tabproof}

\FormulaThmDeltaK[Äquivalenzklasse liegt in der Quotientenmenge]{%
  \EqRel{A,\sim}\dsep x\in A
  \vdash
  \EqClass{x}{\sim}{A}\in \QuotientSet{A}{\sim}
}{EqClassInQuotient}{%
  \DeltaRow{Mengen}{A\dsep x\dsep y}
  \DeltaRow{Relationsoperatoren}{\sim}
  \DeltaRow{Äquivalenzklassen}{\EqClass{x}{\sim}{A}}
}
\begin{tabproofsplit}
  \proofpartindR[Teilmenge des Grundbereichs]{%
    \EqRel{A,\sim}\dsep x\in A
    \vdash \EqClass{x}{\sim}{A}\subseteq A%
  }{%
    \EqRel{A,\sim}\dsep x\in A
    \vdash \EqClass{x}{\sim}{A}\subseteq A%
  }[EqClassSubsetCarrierForQuotient]
    \proofstep{1}{\EqRel{A,\sim}}{\rA}
    \proofstep{2}{x\in A}{\rA}
    \proofstep{3}{y\in \EqClass{x}{\sim}{A}}{\rA}
    \proofstep{1,2,3}{y\in A}{%
      \FormulaRefAuto{EqClassElementInCarrier}{1,2,3}}
    \proofstep{1,2}{\EqClass{x}{\sim}{A}\subseteq A}{%
      \FormulaRefAuto{A\subseteq B \coloneq
        \forall x\,(x\in A \rightarrow x\in B)}{%
        \rUI{\rRI{3,4}}}}
  \closeproofpartindR

  \proofpartindR[Zeuge der Quotientendefinition]{%
    \begin{aligned}[t]
      &\EqRel{A,\sim}\dsep x\in A\vdash{}\\
      &\EqClass{x}{\sim}{A}\in\powerset(A)\land
      \exists z\in A\,
        \EqClass{x}{\sim}{A}=\EqClass{z}{\sim}{A}
    \end{aligned}%
  }{%
    \EqRel{A,\sim}\dsep x\in A\vdash
    \EqClass{x}{\sim}{A}\in\powerset(A)\land
    \exists z\in A\,
      \EqClass{x}{\sim}{A}=\EqClass{z}{\sim}{A}%
  }[EqClassQuotientMembershipWitness]
    \proofstep{1}{\EqRel{A,\sim}}{\rA}
    \proofstep{2}{x\in A}{\rA}
    \proofstep{1,2}{\EqClass{x}{\sim}{A}\subseteq A}{%
      \FormulaRefAuto{EqClassSubsetCarrierForQuotient}{1,2}}
    \proofstep{1,2}{\EqClass{x}{\sim}{A}\in\powerset(A)}{%
      \FormulaRefAuto{x \in \powerset(A)\;\eqvdash\;
        x \subseteq A}{3}}
    \proofstep{}{\EqClass{x}{\sim}{A}=\EqClass{x}{\sim}{A}}{\rII}
    \proofstep{2}{\exists z\in A\,
      \EqClass{x}{\sim}{A}=\EqClass{z}{\sim}{A}}{%
      \rEI{\rAI{2,5}}}
    \proofstep{1,2}{%
      \EqClass{x}{\sim}{A}\in\powerset(A)\land
      \exists z\in A\,
        \EqClass{x}{\sim}{A}=\EqClass{z}{\sim}{A}}{\rAI{4,6}}
  \closeproofpartindR

  \proofpartindR[Einsetzen in die Quotientenmenge]{%
    \EqRel{A,\sim}\dsep x\in A
    \vdash \EqClass{x}{\sim}{A}\in\QuotientSet{A}{\sim}%
  }{%
    \EqRel{A,\sim}\dsep x\in A
    \vdash \EqClass{x}{\sim}{A}\in\QuotientSet{A}{\sim}%
  }[EqClassInQuotientConclusion]
    \proofstep{1}{\EqRel{A,\sim}}{\rA}
    \proofstep{2}{x\in A}{\rA}
    \proofstep{1,2}{%
      \EqClass{x}{\sim}{A}\in\powerset(A)\land
      \exists z\in A\,
        \EqClass{x}{\sim}{A}=\EqClass{z}{\sim}{A}}{%
      \FormulaRefAuto{EqClassQuotientMembershipWitness}{1,2}}
    \proofstep{1,2}{\EqClass{x}{\sim}{A}\in\QuotientSet{A}{\sim}}{%
      \FormulaRefAuto{QuotientSetChar}{1,3}}
  \closeproofpartindR
\end{tabproofsplit}

\FormulaThmDelta[Elemente der Quotientenmenge sind Teilmengen]{%
  \EqRel{A,\sim}\dsep C\in \QuotientSet{A}{\sim}
  \vdash
  C\subseteq A
}{%
  \DeltaRow{Mengen}{A\dsep C}
  \DeltaRow{Relationsoperatoren}{\sim}
}
\begin{tabproof}
  \proofstep{1}{\EqRel{A,\sim}}{\rA}
  \proofstep{2}{C\in \QuotientSet{A}{\sim}}{\rA}
  \proofstep{1,2}{%
    C\in\powerset(A)\land
    \exists x\in A\,C=\EqClass{x}{\sim}{A}}{%
    \FormulaRefAuto{QuotientSetChar}{1,2}}
  \proofstep{1,2}{C\in\powerset(A)}{\rAEa{3}}
  \proofstep{1,2}{C\subseteq A}{%
    \FormulaRefAuto{x \in \powerset(A)\;\eqvdash\; x \subseteq A}{4}}
\end{tabproof}

\FormulaThmDeltaK[Elemente der Quotientenmenge haben einen Repräsentanten]{%
  \EqRel{A,\sim}\dsep C\in \QuotientSet{A}{\sim}
  \vdash
  \exists x\in A\,C=\EqClass{x}{\sim}{A}
}{QuotientRepresentative}{%
  \DeltaRow{Mengen}{A\dsep C\dsep x}
  \DeltaRow{Relationsoperatoren}{\sim}
}
\begin{tabproof}
  \proofstep{1}{\EqRel{A,\sim}}{\rA}
  \proofstep{2}{C\in \QuotientSet{A}{\sim}}{\rA}
  \proofstep{1,2}{%
    C\in\powerset(A)\land
    \exists x\in A\,C=\EqClass{x}{\sim}{A}}{%
    \FormulaRefAuto{QuotientSetChar}{1,2}}
  \proofstep{1,2}{\exists x\in A\,C=\EqClass{x}{\sim}{A}}{\rAEb{3}}
\end{tabproof}


\chapter{Quotientenprojektionen und Quotientenabbildungen}

\section{Quotientenprojektion}

\FormulaDefDeltaK[Quotientenprojektion]{%
  \EqRel{A,\sim}\vdash
  \begin{aligned}[t]
    &\QuotProj{A}{\sim}\colon A\to\QuotientSet{A}{\sim},\\[-2pt]
    &\forall x\in A\,
      \bigl(\QuotProj{A}{\sim}(x)\coloneqq\EqClass{x}{\sim}{A}\bigr)
  \end{aligned}
}{QuotProjDef}{%
  \DeltaRow{Mengen}{A\dsep x}
  \DeltaRow{Funktionen}{\QuotProj{A}{\sim}}
  \DeltaRow{Relationsoperatoren}{\sim}
  \DeltaRow{Quotientenmenge}{\QuotientSet{A}{\sim}}
  \DeltaRow{Äquivalenzklassen}{\EqClass{x}{\sim}{A}}
  \DeltaRow{Neue Symbole}{\QuotProj{A}{\sim}}
}
\begin{tabproof}
  \proofstep{1}{\EqRel{A,\sim}}{\rA}
  \proofstep{2}{x\in A}{\rA}
  \proofstep{1,2}{%
    \EqClass{x}{\sim}{A}\in\QuotientSet{A}{\sim}}{%
    \FormulaRefAuto{EqClassInQuotient}[thm]{1,2}}
\end{tabproof}

\FormulaThmDeltaK[Typisierung der Quotientenprojektion]{%
  \EqRel{A,\sim}\vdash
  \QuotProj{A}{\sim}\colon A\to\QuotientSet{A}{\sim}
}{QuotProjFunction}{%
  \DeltaRow{Mengen}{A}
  \DeltaRow{Funktionen}{\QuotProj{A}{\sim}}
}
\begin{tabproof}
  \proofstep{1}{\EqRel{A,\sim}}{\rA}
  \proofstep{1}{\QuotProj{A}{\sim}\colon A\to\QuotientSet{A}{\sim}}{%
    \FormulaRefAuto{QuotProjDef}[def]{1}}
\end{tabproof}

\FormulaThmDeltaK[Wert der Quotientenprojektion]{%
  \EqRel{A,\sim}\dsep x\in A\vdash
  \QuotProj{A}{\sim}(x)=\EqClass{x}{\sim}{A}
}{QuotProjValue}{%
  \DeltaRow{Mengen}{A\dsep x}
  \DeltaRow{Funktionen}{\QuotProj{A}{\sim}}
  \DeltaRow{Äquivalenzklassen}{\EqClass{x}{\sim}{A}}
}
\begin{tabproof}
  \proofstep{1}{\EqRel{A,\sim}}{\rA}
  \proofstep{2}{x\in A}{\rA}
  \proofstep{1}{\forall u\in A\,
    \bigl(\QuotProj{A}{\sim}(u)=\EqClass{u}{\sim}{A}\bigr)}{%
    \FormulaRefAuto{QuotProjDef}[def]{1}}
  \proofstep{1,2}{\QuotProj{A}{\sim}(x)=\EqClass{x}{\sim}{A}}{%
    \rRE{\rUE{3},2}}
\end{tabproof}

\FormulaThmDeltaK[Surjektivität der Quotientenprojektion]{%
  \EqRel{A,\sim}\vdash
  \QuotProj{A}{\sim}\colon A\sur\QuotientSet{A}{\sim}
}{QuotProjSurjective}{%
  \DeltaRow{Mengen}{A\dsep C\dsep x}
  \DeltaRow{Relationsoperatoren}{\sim}
  \DeltaRow{Funktionen}{\QuotProj{A}{\sim}}
}
\begin{tabproofsplitwide}
  \proofpartwideindR[Jede Quotientenklasse besitzt ein Urbild]{%
    \EqRel{A,\sim}\dsep C\in\QuotientSet{A}{\sim}
    \vdash
    \exists x\in A\,\QuotProj{A}{\sim}(x)=C
  }{%
    \EqRel{A,\sim}\dsep C\in\QuotientSet{A}{\sim}
    \vdash
    \exists x\in A\,\QuotProj{A}{\sim}(x)=C
  }[QuotProjSurjectivePreimage]
    \proofstepwidestar[1]{\EqRel{A,\sim}}{\rA}
    \proofstepwidestar[2]{C\in\QuotientSet{A}{\sim}}{\rA}
    \proofstepwidestar[1,2]{\exists x\in A\,C=\EqClass{x}{\sim}{A}}{%
      \FormulaRefAuto{QuotientRepresentative}{1,2}}
    \proofstepwidestar[4]{x\in A\land C=\EqClass{x}{\sim}{A}}{\rA}
    \proofstepwidestar[4]{x\in A}{\rAEa{4}}
    \proofstepwidestar[4]{C=\EqClass{x}{\sim}{A}}{\rAEb{4}}
    \proofstepwidestar[1,4]{\QuotProj{A}{\sim}(x)=\EqClass{x}{\sim}{A}}{%
      \FormulaRefAuto{QuotProjValue}{1,5}}
    \proofstepwidestar[1,4]{C=\QuotProj{A}{\sim}(x)}{%
      \FormulaRefAuto{a = b,\, c = b \vdash a = c}{6,7}}
    \proofstepwidestar[1,4]{\QuotProj{A}{\sim}(x)=C}{%
      \FormulaRefAuto{a=b\vdash b=a}{8}}
    \proofstepwidestar[1,4]{\exists x\in A\,\QuotProj{A}{\sim}(x)=C}{%
      \rEI{\rAI{5,9}}}
    \proofstepwidestar[1,2]{\exists x\in A\,\QuotProj{A}{\sim}(x)=C}{%
      \rEE{3,4,10}}
  \closeproofpartwideindR

  \pagebreak[3]
  \proofpartwide{Schluss}
    \proofstepwidestar[1]{\EqRel{A,\sim}}{\rA}
    \proofstepwidestar[1]{\QuotProj{A}{\sim}\colon A\to\QuotientSet{A}{\sim}}{%
      \FormulaRefAuto{QuotProjFunction}{1}}
    \proofstepwidestar[3]{C\in\QuotientSet{A}{\sim}}{\rA}
    \proofstepwidestar[1,3]{\exists x\in A\,\QuotProj{A}{\sim}(x)=C}{%
      \FormulaRefAuto{QuotProjSurjectivePreimage}{1,3}}
    \proofstepwidestar[1]{\forall C\in\QuotientSet{A}{\sim}\,
      \exists x\in A\,\QuotProj{A}{\sim}(x)=C}{%
      \rUI{\rRI{3,4}}}
    \proofstepwidestar[1]{\QuotProj{A}{\sim}\colon A\sur\QuotientSet{A}{\sim}}{%
      \FormulaRefAuto{Surjektive Funktion}{2,5}}
  \closeproofpartwide
\end{tabproofsplitwide}

\FormulaThmDeltaK[Gleichheit von Projektionswerten]{%
  \EqRel{A,\sim}\dsep x\in A\dsep y\in A
  \vdash
  \bigl(\QuotProj{A}{\sim}(x)=\QuotProj{A}{\sim}(y)\bigr)
  \leftrightarrow x\sim y
}{QuotProjEqualIffRel}{%
  \DeltaRow{Mengen}{A\dsep x\dsep y}
  \DeltaRow{Relationsoperatoren}{\sim}
  \DeltaRow{Funktionen}{\QuotProj{A}{\sim}}
}
\begin{tabproofsplitwide}
  \proofpartwideindR[Von gleichen Projektionswerten zur Äquivalenz]{%
    \EqRel{A,\sim}\dsep x\in A\dsep y\in A
    \vdash
    \QuotProj{A}{\sim}(x)=\QuotProj{A}{\sim}(y)\rightarrow x\sim y
  }{%
    \EqRel{A,\sim}\dsep x\in A\dsep y\in A
    \vdash
    \QuotProj{A}{\sim}(x)=\QuotProj{A}{\sim}(y)\rightarrow x\sim y
  }[QuotProjEqualIffRelForward]
    \proofstepwidestar[1]{\EqRel{A,\sim}}{\rA}
    \proofstepwidestar[2]{x\in A}{\rA}
    \proofstepwidestar[3]{y\in A}{\rA}
    \proofstepwidestar[4]{\QuotProj{A}{\sim}(x)=\QuotProj{A}{\sim}(y)}{\rA}
    \proofstepwidestar[1,2]{\QuotProj{A}{\sim}(x)=\EqClass{x}{\sim}{A}}{%
      \FormulaRefAuto{QuotProjValue}{1,2}}
    \proofstepwidestar[1,3]{\QuotProj{A}{\sim}(y)=\EqClass{y}{\sim}{A}}{%
      \FormulaRefAuto{QuotProjValue}{1,3}}
    \proofstepwidestar[1,2]{\EqClass{x}{\sim}{A}=\QuotProj{A}{\sim}(x)}{%
      \FormulaRefAuto{a=b\vdash b=a}{5}}
    \proofstepwidestar[1,2,3,4]{\EqClass{x}{\sim}{A}=\QuotProj{A}{\sim}(y)}{%
      \FormulaRefAuto{a = b,\, b = c \vdash a = c}{7,4}}
    \proofstepwidestar[1,2,3,4]{\EqClass{x}{\sim}{A}=\EqClass{y}{\sim}{A}}{%
      \FormulaRefAuto{a = b,\, b = c \vdash a = c}{8,6}}
    \proofstepwidestar[1,2,3,4]{x\sim y}{%
      \FormulaRefAuto{EqClassRelFromEqual}{1,2,3,9}}
    \proofstepwidestar[1,2,3]{%
      \QuotProj{A}{\sim}(x)=\QuotProj{A}{\sim}(y)\rightarrow x\sim y}{%
      \rRI{4,10}}
  \closeproofpartwideindR

  \proofpartwideindR[Von der Äquivalenz zu gleichen Projektionswerten]{%
    \EqRel{A,\sim}\dsep x\in A\dsep y\in A
    \vdash
    x\sim y\rightarrow
    \QuotProj{A}{\sim}(x)=\QuotProj{A}{\sim}(y)
  }{%
    \EqRel{A,\sim}\dsep x\in A\dsep y\in A
    \vdash
    x\sim y\rightarrow
    \QuotProj{A}{\sim}(x)=\QuotProj{A}{\sim}(y)
  }[QuotProjEqualIffRelBackward]
    \proofstepwidestar[1]{\EqRel{A,\sim}}{\rA}
    \proofstepwidestar[2]{x\in A}{\rA}
    \proofstepwidestar[3]{y\in A}{\rA}
    \proofstepwidestar[4]{x\sim y}{\rA}
    \proofstepwidestar[1,2]{\QuotProj{A}{\sim}(x)=\EqClass{x}{\sim}{A}}{%
      \FormulaRefAuto{QuotProjValue}{1,2}}
    \proofstepwidestar[1,3]{\QuotProj{A}{\sim}(y)=\EqClass{y}{\sim}{A}}{%
      \FormulaRefAuto{QuotProjValue}{1,3}}
    \proofstepwidestar[1,2,3,4]{\EqClass{x}{\sim}{A}=\EqClass{y}{\sim}{A}}{%
      \FormulaRefAuto{EqClassEqualFromRel}{1,2,3,4}}
    \proofstepwidestar[1,2,3,4]{\QuotProj{A}{\sim}(x)=\EqClass{y}{\sim}{A}}{%
      \FormulaRefAuto{a = b,\, b = c \vdash a = c}{5,7}}
    \proofstepwidestar[1,3]{\EqClass{y}{\sim}{A}=\QuotProj{A}{\sim}(y)}{%
      \FormulaRefAuto{a=b\vdash b=a}{6}}
    \proofstepwidestar[1,2,3,4]{\QuotProj{A}{\sim}(x)=\QuotProj{A}{\sim}(y)}{%
      \FormulaRefAuto{a = b,\, b = c \vdash a = c}{8,9}}
    \proofstepwidestar[1,2,3]{%
      x\sim y\rightarrow
      \QuotProj{A}{\sim}(x)=\QuotProj{A}{\sim}(y)}{%
      \rRI{4,10}}
  \closeproofpartwideindR

  \proofpartwide{Schluss}
    \proofstepwidestar[1]{\EqRel{A,\sim}}{\rA}
    \proofstepwidestar[2]{x\in A}{\rA}
    \proofstepwidestar[3]{y\in A}{\rA}
    \proofstepwidestar[1,2,3]{%
      \QuotProj{A}{\sim}(x)=\QuotProj{A}{\sim}(y)\rightarrow x\sim y}{%
      \FormulaRefAuto{QuotProjEqualIffRelForward}{1,2,3}}
    \proofstepwidestar[1,2,3]{%
      x\sim y\rightarrow
      \QuotProj{A}{\sim}(x)=\QuotProj{A}{\sim}(y)}{%
      \FormulaRefAuto{QuotProjEqualIffRelBackward}{1,2,3}}
    \proofstepwidestar[1,2,3]{%
      \bigl(\QuotProj{A}{\sim}(x)=\QuotProj{A}{\sim}(y)\bigr)
      \leftrightarrow x\sim y}{\rLRI{4,5}}
  \closeproofpartwide
\end{tabproofsplitwide}

\section{Quotientenbildabbildung einer Teilfamilie}

\FormulaDefDeltaK[Quotientenbildabbildung einer Teilfamilie]{%
  \EqRel{A,\sim}\dsep M\subseteq\powerset(A)\vdash
  \QuotFamMap{A}{\sim}{M}\coloneqq
  \PowMap{\QuotProj{A}{\sim}}\restriction_M
}{QuotFamMapDef}{%
  \DeltaRow{Mengen}{A\dsep M}
  \DeltaRow{Funktionen}{\QuotProj{A}{\sim}\dsep \PowMap{\QuotProj{A}{\sim}}\dsep
  \QuotFamMap{A}{\sim}{M}}
  \DeltaRow{Neue Symbole}{\QuotFamMap{A}{\sim}{M}}
}

\FormulaThmDeltaK[Typisierung der Quotientenbildabbildung]{%
  \EqRel{A,\sim}\dsep M\subseteq\powerset(A)\vdash
  \QuotFamMap{A}{\sim}{M}\colon
  M\to\powerset(\QuotientSet{A}{\sim})
}{QuotFamMapFunction}{%
  \DeltaRow{Mengen}{A\dsep M}
  \DeltaRow{Funktionen}{\QuotProj{A}{\sim}\dsep \QuotFamMap{A}{\sim}{M}}
}
\begin{tabproof}
  \proofstep{1}{\EqRel{A,\sim}}{\rA}
  \proofstep{2}{M\subseteq\powerset(A)}{\rA}
  \proofstep{1}{\QuotProj{A}{\sim}\colon A\to\QuotientSet{A}{\sim}}{%
    \FormulaRefAuto{QuotProjFunction}{1}}
  \proofstep{1}{\PowMap{\QuotProj{A}{\sim}}\colon
    \powerset(A)\to\powerset(\QuotientSet{A}{\sim})}{%
    \FormulaRefAuto{F\colon A\to B \vdash
    \PowMap{F}\colon \powerset(A)\to \powerset(B)}{3}}
  \proofstep{1,2}{\PowMap{\QuotProj{A}{\sim}}\restriction_M
    \colon M\to\powerset(\QuotientSet{A}{\sim})}{%
    \FormulaRefAuto{C\subseteq A \vdash F\restriction_C\colon C\to B}{2,4}}
  \proofstep{1,2}{\QuotFamMap{A}{\sim}{M}=
    \PowMap{\QuotProj{A}{\sim}}\restriction_M}{%
    \FormulaRefAuto{QuotFamMapDef}{1,2}}
  \proofstep{1,2}{\QuotFamMap{A}{\sim}{M}\colon
    M\to\powerset(\QuotientSet{A}{\sim})}{\rIE{6,5}}
\end{tabproof}

\FormulaThmDeltaK[Wert der Quotientenbildabbildung]{%
  \EqRel{A,\sim}\dsep M\subseteq\powerset(A)\dsep X\in M\vdash
  \QuotFamMap{A}{\sim}{M}(X)=\QuotProj{A}{\sim}[X]
}{QuotFamMapValue}{%
  \DeltaRow{Mengen}{A\dsep M\dsep X}
  \DeltaRow{Funktionen}{\QuotProj{A}{\sim}\dsep \QuotFamMap{A}{\sim}{M}}
}
\begin{tabproofsplit}
  \proofpartindR[Auswertung der Einschränkung]{%
    \begin{aligned}[t]
      &\EqRel{A,\sim}\dsep M\subseteq\powerset(A)\dsep X\in M
      \vdash{}\\
      &\QuotFamMap{A}{\sim}{M}(X)
      =\PowMap{\QuotProj{A}{\sim}}(X)
    \end{aligned}%
  }{%
    \EqRel{A,\sim}\dsep M\subseteq\powerset(A)\dsep X\in M
    \vdash
    \QuotFamMap{A}{\sim}{M}(X)
    =\PowMap{\QuotProj{A}{\sim}}(X)%
  }[QuotFamMapValueRestrictedEvaluation]
    \proofstep{1}{\EqRel{A,\sim}}{\rA}
    \proofstep{2}{M\subseteq\powerset(A)}{\rA}
    \proofstep{3}{X\in M}{\rA}
    \proofstep{1}{\QuotProj{A}{\sim}\colon
      A\to\QuotientSet{A}{\sim}}{%
      \FormulaRefAuto{QuotProjFunction}{1}}
    \proofstep{1}{\PowMap{\QuotProj{A}{\sim}}\colon
      \powerset(A)\to\powerset(\QuotientSet{A}{\sim})}{%
      \FormulaRefAuto{F\colon A\to B \vdash
        \PowMap{F}\colon \powerset(A)\to\powerset(B)}{4}}
    \proofstep{1,2}{\QuotFamMap{A}{\sim}{M}=
      \PowMap{\QuotProj{A}{\sim}}\restriction_M}{%
      \FormulaRefAuto{QuotFamMapDef}{1,2}}
    \proofstep{1,2}{\QuotFamMap{A}{\sim}{M}\colon
      M\to\powerset(\QuotientSet{A}{\sim})}{%
      \FormulaRefAuto{QuotFamMapFunction}[thm]{1,2}}
    \proofstep{1,2}{\PowMap{\QuotProj{A}{\sim}}\restriction_M\colon
      M\to\powerset(\QuotientSet{A}{\sim})}{%
      \FormulaRefAuto{C\subseteq A \vdash F\restriction_C\colon C\to B}{2,5}}
    \proofstep{1,2,3}{\QuotFamMap{A}{\sim}{M}(X)=
      \bigl(\PowMap{\QuotProj{A}{\sim}}\restriction_M\bigr)(X)}{%
      \FormulaRefAuto{FunctionEqualityEvaluation}[thm]{7,8,3,6}}
    \proofstep{1,2,3}{%
      \bigl(\PowMap{\QuotProj{A}{\sim}}\restriction_M\bigr)(X)
      =\PowMap{\QuotProj{A}{\sim}}(X)}{%
      \FormulaRefAuto{C\subseteq A \dsep x\in C
        \vdash F\restriction_C(x)=F(x)}{2,3,5}}
    \proofstep{1,2,3}{\QuotFamMap{A}{\sim}{M}(X)
      =\PowMap{\QuotProj{A}{\sim}}(X)}{%
      \FormulaRefAuto{a = b,\, b = c \vdash a = c}{9,10}}
  \closeproofpartindR

  \proofpartindR[Potenzmengenabbildung als Bild]{%
    \begin{aligned}[t]
      &\EqRel{A,\sim}\dsep M\subseteq\powerset(A)\dsep X\in M
      \vdash{}\\
      &\PowMap{\QuotProj{A}{\sim}}(X)=\QuotProj{A}{\sim}[X]
    \end{aligned}%
  }{%
    \EqRel{A,\sim}\dsep M\subseteq\powerset(A)\dsep X\in M
    \vdash
    \PowMap{\QuotProj{A}{\sim}}(X)=\QuotProj{A}{\sim}[X]%
  }[QuotFamMapValuePowerImage]
    \proofstep{1}{\EqRel{A,\sim}}{\rA}
    \proofstep{2}{M\subseteq\powerset(A)}{\rA}
    \proofstep{3}{X\in M}{\rA}
    \proofstep{1}{\QuotProj{A}{\sim}\colon
      A\to\QuotientSet{A}{\sim}}{%
      \FormulaRefAuto{QuotProjFunction}{1}}
    \proofstep{2,3}{X\in\powerset(A)}{%
      \FormulaRefAuto{A\subseteq B,\,x\in A\vdash x\in B}{2,3}}
    \proofstep{1,2,3}{\PowMap{\QuotProj{A}{\sim}}(X)=
      \QuotProj{A}{\sim}[X]}{%
      \FormulaRefAuto{F\colon A\to B\dsep X\in\powerset(A)\vdash
        \PowMap{F}(X)=F[X]}{4,5}}
  \closeproofpartindR

  \proofpartindR[Zusammenführung]{%
    \EqRel{A,\sim}\dsep M\subseteq\powerset(A)\dsep X\in M
    \vdash
    \QuotFamMap{A}{\sim}{M}(X)=\QuotProj{A}{\sim}[X]%
  }{%
    \EqRel{A,\sim}\dsep M\subseteq\powerset(A)\dsep X\in M
    \vdash
    \QuotFamMap{A}{\sim}{M}(X)=\QuotProj{A}{\sim}[X]%
  }[QuotFamMapValueConclusion]
    \proofstep{1}{\EqRel{A,\sim}}{\rA}
    \proofstep{2}{M\subseteq\powerset(A)}{\rA}
    \proofstep{3}{X\in M}{\rA}
    \proofstep{1,2,3}{\QuotFamMap{A}{\sim}{M}(X)
      =\PowMap{\QuotProj{A}{\sim}}(X)}{%
      \FormulaRefAuto{QuotFamMapValueRestrictedEvaluation}{1,2,3}}
    \proofstep{1,2,3}{\PowMap{\QuotProj{A}{\sim}}(X)
      =\QuotProj{A}{\sim}[X]}{%
      \FormulaRefAuto{QuotFamMapValuePowerImage}{1,2,3}}
    \proofstep{1,2,3}{\QuotFamMap{A}{\sim}{M}(X)
      =\QuotProj{A}{\sim}[X]}{%
      \FormulaRefAuto{a = b,\, b = c \vdash a = c}{4,5}}
  \closeproofpartindR
\end{tabproofsplit}

\subsection{Charakterisierung der Quotientenbildabbildung}

\FormulaThmDeltaK[Quotientenprojektionsbilder liegen in der Quotientenmenge]{%
  \EqRel{A,\sim}\dsep M\subseteq\powerset(A)\dsep X\in M
  \dsep C\in\QuotProj{A}{\sim}[X]
  \vdash
  C\in\QuotientSet{A}{\sim}
}{QuotProjImageInQuotient}{%
  \DeltaRow{Mengen}{A\dsep M\dsep X\dsep C\dsep a}
  \DeltaRow{Funktionen}{\QuotProj{A}{\sim}}
  \DeltaRow{Quotientenmenge}{\QuotientSet{A}{\sim}}
}
\begin{tabproof}
  \proofstep{1}{\EqRel{A,\sim}}{\rA}
  \proofstep{2}{M\subseteq\powerset(A)}{\rA}
  \proofstep{3}{X\in M}{\rA}
  \proofstep{4}{C\in\QuotProj{A}{\sim}[X]}{\rA}
  \proofstep{2,3}{X\in\powerset(A)}{%
    \FormulaRefAuto{A\subseteq B,\, x\in A \vdash x\in B}{2,3}}
  \proofstep{2,3}{X\subseteq A}{%
    \FormulaRefAuto{x \in \powerset(A)\;\eqvdash\; x \subseteq A}{5}}
  \proofstep{1}{\QuotProj{A}{\sim}\colon
    A\to\QuotientSet{A}{\sim}}{%
    \FormulaRefAuto{QuotProjFunction}{1}}
  \proofstep{1,2,3}{\QuotProj{A}{\sim}[X]\subseteq
    \QuotientSet{A}{\sim}}{%
    \FormulaRefAuto{C\subseteq A \dsep F\colon A\to B \vdash F[C]\subseteq B}{6,7}}
  \proofstep{1,2,3,4}{C\in\QuotientSet{A}{\sim}}{%
    \FormulaRefAuto{A\subseteq B,\, x\in A \vdash x\in B}{8,4}}
\end{tabproof}

\FormulaThmDeltaK[Quotientenprojektionsbilder liefern Klassenzeugen]{%
  \EqRel{A,\sim}\dsep M\subseteq\powerset(A)\dsep X\in M
  \dsep C\in\QuotProj{A}{\sim}[X]
  \vdash
  \exists a\in X\,C=\EqClass{a}{\sim}{A}
}{QuotProjImageClassWitness}{%
  \DeltaRow{Mengen}{A\dsep M\dsep X\dsep C\dsep a}
  \DeltaRow{Funktionen}{\QuotProj{A}{\sim}}
  \DeltaRow{Quotientenmenge}{\QuotientSet{A}{\sim}}
}
\begin{tabproofsplitwide}
  \proofpartwideindR[Urbildzeuge des Projektionsbildes]{%
    \EqRel{A,\sim}\dsep M\subseteq\powerset(A)\dsep X\in M
    \dsep C\in\QuotProj{A}{\sim}[X]
    \vdash
    \exists a\in X\,C=\QuotProj{A}{\sim}(a)
  }{%
    \EqRel{A,\sim}\dsep M\subseteq\powerset(A)\dsep X\in M
    \dsep C\in\QuotProj{A}{\sim}[X]
    \vdash
    \exists a\in X\,C=\QuotProj{A}{\sim}(a)
  }[QuotProjImageValueWitness]
    \proofstepwidestar[1]{\EqRel{A,\sim}}{\rA}
    \proofstepwidestar[2]{M\subseteq\powerset(A)}{\rA}
    \proofstepwidestar[3]{X\in M}{\rA}
    \proofstepwidestar[4]{C\in\QuotProj{A}{\sim}[X]}{\rA}
    \proofstepwidestar[2,3]{X\in\powerset(A)}{%
      \FormulaRefAuto{A\subseteq B,\, x\in A \vdash x\in B}{2,3}}
    \proofstepwidestar[2,3]{X\subseteq A}{%
      \FormulaRefAuto{x \in \powerset(A)\;\eqvdash\; x \subseteq A}{5}}
    \proofstepwidestar[1]{\QuotProj{A}{\sim}\colon
      A\to\QuotientSet{A}{\sim}}{%
      \FormulaRefAuto{QuotProjFunction}{1}}
    \proofstepwidestar[1,2,3,4]{\exists a\in X\,
      C=\QuotProj{A}{\sim}(a)}{%
      \FormulaRefAuto{C\subseteq A\dsep F\colon A\to B\dsep
      y\in F[C]\vdash \exists x\in C\,y=F(x)}{6,7,4}}
  \closeproofpartwideindR

  \proofpartwideindR[Vom Urbildzeugen zum Klassenzeugen]{%
    \EqRel{A,\sim}\dsep M\subseteq\powerset(A)\dsep X\in M
    \dsep C\in\QuotProj{A}{\sim}[X]
    \vdash
    \exists a\in X\,C=\EqClass{a}{\sim}{A}
  }{%
    \EqRel{A,\sim}\dsep M\subseteq\powerset(A)\dsep X\in M
    \dsep C\in\QuotProj{A}{\sim}[X]
    \vdash
    \exists a\in X\,C=\EqClass{a}{\sim}{A}
  }[QuotProjImageClassWitnessConclusion]
    \proofstepwidestar[1]{\EqRel{A,\sim}}{\rA}
    \proofstepwidestar[2]{M\subseteq\powerset(A)}{\rA}
    \proofstepwidestar[3]{X\in M}{\rA}
    \proofstepwidestar[4]{C\in\QuotProj{A}{\sim}[X]}{\rA}
    \proofstepwidestar[1,2,3,4]{\exists a\in X\,
      C=\QuotProj{A}{\sim}(a)}{%
      \FormulaRefAuto{QuotProjImageValueWitness}{1,2,3,4}}
    \proofstepwidestar[6]{a\in X\land C=\QuotProj{A}{\sim}(a)}{\rA}
    \proofstepwidestar[6]{a\in X}{\rAEa{6}}
    \proofstepwidestar[2,3,6]{X\subseteq A}{%
      \FormulaRefAuto{x \in \powerset(A)\;\eqvdash\; x \subseteq A}{%
        \FormulaRefAuto{A\subseteq B,\, x\in A \vdash x\in B}{2,3}}}
    \proofstepwidestar[2,3,6]{a\in A}{%
      \FormulaRefAuto{A\subseteq B,\, x\in A \vdash x\in B}{8,7}}
    \proofstepwidestar[1,6]{\QuotProj{A}{\sim}(a)=\EqClass{a}{\sim}{A}}{%
      \FormulaRefAuto{QuotProjValue}{1,9}}
    \proofstepwidestar[1,6]{C=\EqClass{a}{\sim}{A}}{%
      \FormulaRefAuto{a = b,\, b = c \vdash a = c}{\rAEb{6},10}}
    \proofstepwidestar[1,6]{a\in X\land C=\EqClass{a}{\sim}{A}}{\rAI{7,11}}
    \proofstepwidestar[1,6]{\exists a\in X\,C=\EqClass{a}{\sim}{A}}{\rEI{12}}
    \proofstepwidestar[1,2,3,4]{\exists a\in X\,C=\EqClass{a}{\sim}{A}}{%
      \rEE{5,6,13}}
  \closeproofpartwideindR
\end{tabproofsplitwide}

\FormulaThmDeltaK[Quotientenprojektionsbilder liefern Klassen]{%
  \begin{aligned}[t]
    &\EqRel{A,\sim}\dsep M\subseteq\powerset(A)\dsep X\in M
      \dsep C\in\QuotProj{A}{\sim}[X]\vdash{}\\
    &C\in\QuotientSet{A}{\sim}\land
      \exists a\in X\,C=\EqClass{a}{\sim}{A}
  \end{aligned}
}{QuotProjImageToEqClass}{%
  \DeltaRow{Mengen}{A\dsep M\dsep X\dsep C\dsep a}
  \DeltaRow{Funktionen}{\QuotProj{A}{\sim}}
  \DeltaRow{Quotientenmenge}{\QuotientSet{A}{\sim}}
}
\begin{tabproof}
  \proofstep{1}{\EqRel{A,\sim}}{\rA}
  \proofstep{2}{M\subseteq\powerset(A)}{\rA}
  \proofstep{3}{X\in M}{\rA}
  \proofstep{4}{C\in\QuotProj{A}{\sim}[X]}{\rA}
  \proofstep{1,2,3,4}{C\in\QuotientSet{A}{\sim}}{%
    \FormulaRefAuto{QuotProjImageInQuotient}{1,2,3,4}}
  \proofstep{1,2,3,4}{\exists a\in X\,C=\EqClass{a}{\sim}{A}}{%
    \FormulaRefAuto{QuotProjImageClassWitness}{1,2,3,4}}
  \proofstep{1,2,3,4}{C\in\QuotientSet{A}{\sim}\land
    \exists a\in X\,C=\EqClass{a}{\sim}{A}}{\rAI{5,6}}
\end{tabproof}

\FormulaThmDeltaK[Klassen liefern Quotientenprojektionsbilder]{%
  \EqRel{A,\sim}\dsep M\subseteq\powerset(A)\dsep X\in M
  \dsep \exists a\in X\,C=\EqClass{a}{\sim}{A}
  \vdash
  C\in\QuotProj{A}{\sim}[X]
}{EqClassToQuotProjImage}{%
  \DeltaRow{Mengen}{A\dsep M\dsep X\dsep C\dsep a}
  \DeltaRow{Funktionen}{\QuotProj{A}{\sim}}
  \DeltaRow{Quotientenmenge}{\QuotientSet{A}{\sim}}
}
\begin{tabproofsplitwide}
  \proofpartwideindR[Ein Klassenzeuge liefert einen Projektionswert]{%
    \EqRel{A,\sim}\dsep M\subseteq\powerset(A)\dsep X\in M
    \dsep a\in X\land C=\EqClass{a}{\sim}{A}
    \vdash
    C\in\QuotProj{A}{\sim}[X]
  }{%
    \EqRel{A,\sim}\dsep M\subseteq\powerset(A)\dsep X\in M
    \dsep a\in X\land C=\EqClass{a}{\sim}{A}
    \vdash
    C\in\QuotProj{A}{\sim}[X]
  }[EqClassWitnessToQuotProjImage]
    \proofstepwidestar[1]{\EqRel{A,\sim}}{\rA}
    \proofstepwidestar[2]{M\subseteq\powerset(A)}{\rA}
    \proofstepwidestar[3]{X\in M}{\rA}
    \proofstepwidestar[4]{a\in X\land C=\EqClass{a}{\sim}{A}}{\rA}
    \proofstepwidestar[2,3]{X\in\powerset(A)}{%
      \FormulaRefAuto{A\subseteq B,\, x\in A \vdash x\in B}{2,3}}
    \proofstepwidestar[2,3]{X\subseteq A}{%
      \FormulaRefAuto{x \in \powerset(A)\;\eqvdash\; x \subseteq A}{5}}
    \proofstepwidestar[1]{\QuotProj{A}{\sim}\colon
      A\to\QuotientSet{A}{\sim}}{%
      \FormulaRefAuto{QuotProjFunction}{1}}
    \proofstepwidestar[4]{a\in X}{\rAEa{4}}
    \proofstepwidestar[2,3,4]{a\in A}{%
      \FormulaRefAuto{A\subseteq B,\, x\in A \vdash x\in B}{6,8}}
    \proofstepwidestar[1,4]{\QuotProj{A}{\sim}(a)=\EqClass{a}{\sim}{A}}{%
      \FormulaRefAuto{QuotProjValue}{1,9}}
    \proofstepwidestar[1,4]{\EqClass{a}{\sim}{A}=\QuotProj{A}{\sim}(a)}{%
      \FormulaRefAuto{a=b\vdash b=a}{10}}
    \proofstepwidestar[1,4]{C=\QuotProj{A}{\sim}(a)}{%
      \FormulaRefAuto{a = b,\, b = c \vdash a = c}{\rAEb{4},11}}
    \proofstepwidestar[1,4]{a\in X\land C=\QuotProj{A}{\sim}(a)}{%
      \rAI{8,12}}
    \proofstepwidestar[1,4]{\exists a\in X\,
      C=\QuotProj{A}{\sim}(a)}{\rEI{13}}
    \proofstepwidestar[1,2,3,4]{C\in\QuotProj{A}{\sim}[X]}{%
      \FormulaRefAuto{C\subseteq A\dsep F\colon A\to B\dsep
      \exists x\in C\,y=F(x)\vdash y\in F[C]}{6,7,14}}
  \closeproofpartwideindR

  \proofpartwideindR[Elimination des Klassenzeugen]{%
    \EqRel{A,\sim}\dsep M\subseteq\powerset(A)\dsep X\in M
    \dsep \exists a\in X\,C=\EqClass{a}{\sim}{A}
    \vdash
    C\in\QuotProj{A}{\sim}[X]
  }{%
    \EqRel{A,\sim}\dsep M\subseteq\powerset(A)\dsep X\in M
    \dsep \exists a\in X\,C=\EqClass{a}{\sim}{A}
    \vdash
    C\in\QuotProj{A}{\sim}[X]
  }[EqClassToQuotProjImageConclusion]
    \proofstepwidestar[1]{\EqRel{A,\sim}}{\rA}
    \proofstepwidestar[2]{M\subseteq\powerset(A)}{\rA}
    \proofstepwidestar[3]{X\in M}{\rA}
    \proofstepwidestar[4]{\exists a\in X\,C=\EqClass{a}{\sim}{A}}{\rA}
    \proofstepwidestar[5]{a\in X\land C=\EqClass{a}{\sim}{A}}{\rA}
    \proofstepwidestar[1,2,3,5]{C\in\QuotProj{A}{\sim}[X]}{%
      \FormulaRefAuto{EqClassWitnessToQuotProjImage}{1,2,3,5}}
    \proofstepwidestar[1,2,3,4]{C\in\QuotProj{A}{\sim}[X]}{%
      \rEE{4,5,6}}
  \closeproofpartwideindR
\end{tabproofsplitwide}

\FormulaThmDeltaK[Charakterisierung der Quotientenbildabbildung]{%
  \EqRel{A,\sim}\dsep M\subseteq\powerset(A)\dsep X\in M\vdash
  \begin{aligned}[t]
    C\in\QuotFamMap{A}{\sim}{M}(X)
    \leftrightarrow{}&
    C\in\QuotientSet{A}{\sim}\land{}\\
    &\exists a\in X\,C=\EqClass{a}{\sim}{A}
  \end{aligned}
}{QuotFamMapChar}{%
  \DeltaRow{Mengen}{A\dsep M\dsep X\dsep C\dsep a}
  \DeltaRow{Funktionen}{\QuotProj{A}{\sim}\dsep \QuotFamMap{A}{\sim}{M}}
  \DeltaRow{Quotientenmenge}{\QuotientSet{A}{\sim}}
}
\begin{tabproofsplitwide}
  \proofpartwideindR[\(\rightarrow\)]{%
    \begin{aligned}[t]
      &\EqRel{A,\sim}\dsep M\subseteq\powerset(A)\dsep X\in M
        \vdash{}\\
      &C\in\QuotFamMap{A}{\sim}{M}(X)\rightarrow{}\\[-2pt]
      &\qquad\bigl(C\in\QuotientSet{A}{\sim}\land
        \exists a\in X\,C=\EqClass{a}{\sim}{A}\bigr)
    \end{aligned}
  }{%
    \begin{aligned}[t]
      &\EqRel{A,\sim}\dsep M\subseteq\powerset(A)\dsep X\in M
        \vdash{}\\
      &C\in\QuotFamMap{A}{\sim}{M}(X)\rightarrow{}\\[-2pt]
      &\qquad\bigl(C\in\QuotientSet{A}{\sim}\land
        \exists a\in X\,C=\EqClass{a}{\sim}{A}\bigr)
    \end{aligned}
  }[QuotFamMapCharForward]
    \proofstepwidestar[1]{\EqRel{A,\sim}}{\rA}
    \proofstepwidestar[2]{M\subseteq\powerset(A)}{\rA}
    \proofstepwidestar[3]{X\in M}{\rA}
    \proofstepwidestar[4]{C\in\QuotFamMap{A}{\sim}{M}(X)}{\rA}
    \proofstepwidestar[1,2,3]{\QuotFamMap{A}{\sim}{M}(X)=
      \QuotProj{A}{\sim}[X]}{%
      \FormulaRefAuto{QuotFamMapValue}{1,2,3}}
    \proofstepwidestar[1,2,3,4]{C\in\QuotProj{A}{\sim}[X]}{\rIE{5,4}}
    \proofstepwidestar[1,2,3,4]{C\in\QuotientSet{A}{\sim}\land
      \exists a\in X\,C=\EqClass{a}{\sim}{A}}{%
      \FormulaRefAuto{QuotProjImageToEqClass}{1,2,3,6}}
    \proofstepwidestar[1,2,3]{C\in\QuotFamMap{A}{\sim}{M}(X)
      \rightarrow
      \bigl(C\in\QuotientSet{A}{\sim}\land
      \exists a\in X\,C=\EqClass{a}{\sim}{A}\bigr)}{\rRI{4,7}}
  \closeproofpartwideindR

  \proofpartwideindR[\(\leftarrow\)]{%
    \begin{aligned}[t]
      &\EqRel{A,\sim}\dsep M\subseteq\powerset(A)\dsep X\in M
        \vdash{}\\
      &\bigl(C\in\QuotientSet{A}{\sim}\land
        \exists a\in X\,C=\EqClass{a}{\sim}{A}\bigr)\rightarrow{}\\[-2pt]
      &\qquad C\in\QuotFamMap{A}{\sim}{M}(X)
    \end{aligned}
  }{%
    \begin{aligned}[t]
      &\EqRel{A,\sim}\dsep M\subseteq\powerset(A)\dsep X\in M
        \vdash{}\\
      &\bigl(C\in\QuotientSet{A}{\sim}\land
        \exists a\in X\,C=\EqClass{a}{\sim}{A}\bigr)\rightarrow{}\\[-2pt]
      &\qquad C\in\QuotFamMap{A}{\sim}{M}(X)
    \end{aligned}
  }[QuotFamMapCharBackward]
    \proofstepwidestar[1]{\EqRel{A,\sim}}{\rA}
    \proofstepwidestar[2]{M\subseteq\powerset(A)}{\rA}
    \proofstepwidestar[3]{X\in M}{\rA}
    \proofstepwidestar[4]{%
      \begin{aligned}[t]
        &C\in\QuotientSet{A}{\sim}\land{}\\[-2pt]
        &\exists a\in X\,C=\EqClass{a}{\sim}{A}
      \end{aligned}}{\rA}
    \proofstepwidestar[4]{\exists a\in X\,C=\EqClass{a}{\sim}{A}}{\rAEb{4}}
    \proofstepwidestar[1,2,3,4]{C\in\QuotProj{A}{\sim}[X]}{%
      \FormulaRefAuto{EqClassToQuotProjImage}{1,2,3,5}}
    \proofstepwidestar[1,2,3]{\QuotProj{A}{\sim}[X]=
      \QuotFamMap{A}{\sim}{M}(X)}{%
      \FormulaRefAuto{a=b\vdash b=a}{%
      \FormulaRefAuto{QuotFamMapValue}{1,2,3}}}
    \proofstepwidestar[1,2,3,4]{C\in\QuotFamMap{A}{\sim}{M}(X)}{%
      \rIE{7,6}}
    \proofstepwidestar[1,2,3]{%
      \begin{aligned}[t]
        &\bigl(C\in\QuotientSet{A}{\sim}\land
          \exists a\in X\,C=\EqClass{a}{\sim}{A}\bigr)\rightarrow{}\\[-2pt]
        &\qquad C\in\QuotFamMap{A}{\sim}{M}(X)
      \end{aligned}}{\rRI{4,8}}
  \closeproofpartwideindR

  \proofpartwide{Schluss}
    \proofstepwidestar[1]{\EqRel{A,\sim}}{\rA}
    \proofstepwidestar[2]{M\subseteq\powerset(A)}{\rA}
    \proofstepwidestar[3]{X\in M}{\rA}
    \proofstepwidestar[1,2,3]{C\in\QuotFamMap{A}{\sim}{M}(X)
      \rightarrow
      \bigl(C\in\QuotientSet{A}{\sim}\land
      \exists a\in X\,C=\EqClass{a}{\sim}{A}\bigr)}{%
      \FormulaRefAuto{QuotFamMapCharForward}{1,2,3}}
    \proofstepwidestar[1,2,3]{%
      \bigl(C\in\QuotientSet{A}{\sim}\land
      \exists a\in X\,C=\EqClass{a}{\sim}{A}\bigr)
      \rightarrow C\in\QuotFamMap{A}{\sim}{M}(X)}{%
      \FormulaRefAuto{QuotFamMapCharBackward}{1,2,3}}
    \proofstepwidestar[1,2,3]{C\in\QuotFamMap{A}{\sim}{M}(X)
      \leftrightarrow
      \bigl(C\in\QuotientSet{A}{\sim}\land
      \exists a\in X\,C=\EqClass{a}{\sim}{A}\bigr)}{\rLRI{4,5}}
  \closeproofpartwide
\end{tabproofsplitwide}

\FormulaThmDeltaK[Quotientenbildabbildungen erhalten Vereinigungen]{%
  \begin{aligned}[t]
    &\EqRel{A,\sim}\dsep M\subseteq\powerset(A)\dsep
      X\in M\dsep Y\in M\dsep X\cup Y\in M\\
    &\vdash
      \QuotFamMap{A}{\sim}{M}(X\cup Y)=
      \QuotFamMap{A}{\sim}{M}(X)\cup\QuotFamMap{A}{\sim}{M}(Y)
  \end{aligned}
}{QuotFamMapPreservesUnion}{%
  \DeltaRow{Mengen}{A\dsep M\dsep X\dsep Y}
  \DeltaRow{Relationsoperatoren}{\sim}
  \DeltaRow{Funktionen}{\QuotProj{A}{\sim}\dsep\QuotFamMap{A}{\sim}{M}}
}
\begin{tabproofsplitwide}
  \proofpartwideindR[Die Quotientenprojektion erhält Vereinigungen von Teilmengen]{%
    \EqRel{A,\sim}\dsep M\subseteq\powerset(A)\dsep X\in M\dsep Y\in M
    \vdash
    \QuotProj{A}{\sim}[X\cup Y]=
    \QuotProj{A}{\sim}[X]\cup\QuotProj{A}{\sim}[Y]
  }{%
    \EqRel{A,\sim}\dsep M\subseteq\powerset(A)\dsep X\in M\dsep Y\in M
    \vdash
    \QuotProj{A}{\sim}[X\cup Y]=
    \QuotProj{A}{\sim}[X]\cup\QuotProj{A}{\sim}[Y]
  }[QuotProjImagePreservesUnion]
    \proofstepwidestar[1]{\EqRel{A,\sim}}{\rA}
    \proofstepwidestar[2]{M\subseteq\powerset(A)}{\rA}
    \proofstepwidestar[3]{X\in M}{\rA}
    \proofstepwidestar[4]{Y\in M}{\rA}
    \proofstepwidestar[1]{\QuotProj{A}{\sim}\colon
      A\to\QuotientSet{A}{\sim}}{%
      \FormulaRefAuto{QuotProjFunction}{1}}
    \proofstepwidestar[2,3]{X\in\powerset(A)}{%
      \FormulaRefAuto{A\subseteq B,\,x\in A\vdash x\in B}{2,3}}
    \proofstepwidestar[2,3]{X\subseteq A}{%
      \FormulaRefAuto{x\in\powerset(A)\eqvdash x\subseteq A}{6}}
    \proofstepwidestar[2,4]{Y\in\powerset(A)}{%
      \FormulaRefAuto{A\subseteq B,\,x\in A\vdash x\in B}{2,4}}
    \proofstepwidestar[2,4]{Y\subseteq A}{%
      \FormulaRefAuto{x\in\powerset(A)\eqvdash x\subseteq A}{8}}
    \proofstepwidestar[1,2,3,4]{%
      \QuotProj{A}{\sim}[X\cup Y]=
      \QuotProj{A}{\sim}[X]\cup\QuotProj{A}{\sim}[Y]}{%
      \FormulaRefAuto{F\colon M\to N \dsep A\subseteq M \dsep
      B\subseteq M \vdash F[A\cup B]=F[A]\cup F[B]}{5,7,9}}
  \closeproofpartwideindR

  \proofpartwideindR[Auswertung auf den beiden Familiengliedern]{%
    \EqRel{A,\sim}\dsep M\subseteq\powerset(A)\dsep X\in M\dsep Y\in M
    \vdash
    \QuotProj{A}{\sim}[X]\cup\QuotProj{A}{\sim}[Y]=
    \QuotFamMap{A}{\sim}{M}(X)\cup\QuotFamMap{A}{\sim}{M}(Y)
  }{%
    \EqRel{A,\sim}\dsep M\subseteq\powerset(A)\dsep X\in M\dsep Y\in M
    \vdash
    \QuotProj{A}{\sim}[X]\cup\QuotProj{A}{\sim}[Y]=
    \QuotFamMap{A}{\sim}{M}(X)\cup\QuotFamMap{A}{\sim}{M}(Y)
  }[QuotFamMapUnionValue]
    \proofstepwidestar[1]{\EqRel{A,\sim}}{\rA}
    \proofstepwidestar[2]{M\subseteq\powerset(A)}{\rA}
    \proofstepwidestar[3]{X\in M}{\rA}
    \proofstepwidestar[4]{Y\in M}{\rA}
    \proofstepwidestar[1,2,3]{%
      \QuotFamMap{A}{\sim}{M}(X)=\QuotProj{A}{\sim}[X]}{%
      \FormulaRefAuto{QuotFamMapValue}{1,2,3}}
    \proofstepwidestar[1,2,4]{%
      \QuotFamMap{A}{\sim}{M}(Y)=\QuotProj{A}{\sim}[Y]}{%
      \FormulaRefAuto{QuotFamMapValue}{1,2,4}}
    \proofstepwidestar[1,2,3]{%
      \QuotProj{A}{\sim}[X]=\QuotFamMap{A}{\sim}{M}(X)}{%
      \FormulaRefAuto{a=b\vdash b=a}{5}}
    \proofstepwidestar[1,2,4]{%
      \QuotProj{A}{\sim}[Y]=\QuotFamMap{A}{\sim}{M}(Y)}{%
      \FormulaRefAuto{a=b\vdash b=a}{6}}
    \proofstepwidestar[1,2,3]{%
      \QuotProj{A}{\sim}[X]\cup\QuotProj{A}{\sim}[Y]=
      \QuotFamMap{A}{\sim}{M}(X)\cup\QuotProj{A}{\sim}[Y]}{%
      \FormulaRefAuto{A=B\vdash A\cup C = B\cup C}{7}}
    \proofstepwidestar[1,2,3,4]{%
      \QuotProj{A}{\sim}[X]\cup\QuotProj{A}{\sim}[Y]=
      \QuotFamMap{A}{\sim}{M}(X)\cup\QuotFamMap{A}{\sim}{M}(Y)}{%
      \FormulaRefAuto{a = b,\, b = c \vdash a = c}{%
        9,\FormulaRefAuto{A=B\vdash C\cup A = C\cup B}{8}}}
  \closeproofpartwideindR

  \proofpartwide{Schluss}
    \proofstepwidestar[1]{\EqRel{A,\sim}}{\rA}
    \proofstepwidestar[2]{M\subseteq\powerset(A)}{\rA}
    \proofstepwidestar[3]{X\in M}{\rA}
    \proofstepwidestar[4]{Y\in M}{\rA}
    \proofstepwidestar[5]{X\cup Y\in M}{\rA}
    \proofstepwidestar[1,2,5]{%
      \QuotFamMap{A}{\sim}{M}(X\cup Y)=\QuotProj{A}{\sim}[X\cup Y]}{%
      \FormulaRefAuto{QuotFamMapValue}{1,2,5}}
    \proofstepwidestar[1,2,3,4]{%
      \QuotProj{A}{\sim}[X\cup Y]=
      \QuotProj{A}{\sim}[X]\cup\QuotProj{A}{\sim}[Y]}{%
      \FormulaRefAuto{QuotProjImagePreservesUnion}{1,2,3,4}}
    \proofstepwidestar[1,2,3,4]{%
      \QuotProj{A}{\sim}[X]\cup\QuotProj{A}{\sim}[Y]=
      \QuotFamMap{A}{\sim}{M}(X)\cup\QuotFamMap{A}{\sim}{M}(Y)}{%
      \FormulaRefAuto{QuotFamMapUnionValue}{1,2,3,4}}
    \proofstepwidestar[1,2,3,4,5]{%
      \QuotFamMap{A}{\sim}{M}(X\cup Y)=
      \QuotFamMap{A}{\sim}{M}(X)\cup\QuotFamMap{A}{\sim}{M}(Y)}{%
      \FormulaRefAuto{a = b,\, b = c \vdash a = c}{%
        6,\FormulaRefAuto{a = b,\, b = c \vdash a = c}{7,8}}}
  \closeproofpartwide
\end{tabproofsplitwide}

\section{Beispiel: Ununterscheidbarkeitsquotient einer Mengenfamilie}

Die allgemeine Quotiententheorie wird nun auf die bereits eingeführte
Ununterscheidbarkeitsrelation einer Mengenfamilie spezialisiert. Dabei
unterscheiden wir ausdrücklich die punktweise Quotientenprojektion von der
durch sie induzierten Quotientenbildabbildung auf der Familie. Die folgenden
Symbole sind nur Abkürzungen für die allgemeinen Konstruktionen dieses
Kapitels.

\begin{notation*}
Für eine Mengenfamilie \(M\) schreiben wir
\[
  \begin{aligned}
    \OccClass{M}{a}
      &\coloneqq \EqClass{a}{\SameOccRel{M}}{\bigcup M},\\
    \OccQuotCarrier{M}
      &\coloneqq \QuotientSet{\bigcup M}{\SameOccRel{M}},\\
    \OccQuotProj{M}
      &\coloneqq \QuotProj{\bigcup M}{\SameOccRel{M}},\\
    \OccQuotMap{M}
      &\coloneqq \QuotFamMap{\bigcup M}{\SameOccRel{M}}{M},\\
    \OccQuotFam{M}
      &\coloneqq \OccQuotMap{M}[M].
  \end{aligned}
\]
\end{notation*}

\subsection{Die punktweise Ununterscheidbarkeitsprojektion}

Die Abbildung \(\OccQuotProj{M}\) schickt jedes Element des Trägers auf seine
Ununterscheidbarkeitsklasse. Sie ist damit die unmittelbare Spezialisierung
der allgemeinen Quotientenprojektion.

\FormulaThmDeltaK[Surjektivität der Ununterscheidbarkeitsprojektion]{%
  \OccQuotProj{M}\colon
  \bigcup M\sur\OccQuotCarrier{M}
}{OccQuotProjSurjective}{%
  \DeltaRow{Mengen}{M}
  \DeltaRow{Relationsoperatoren}{\SameOccRel{M}}
  \DeltaRow{Funktionen}{\OccQuotProj{M}}
}
\begin{tabproof}
  \proofstep{}{\EqRel{\bigcup M,\SameOccRel{M}}}{%
    \FormulaRefAuto{SameOccRelEqRel}}
  \proofstep{}{\OccQuotProj{M}\colon
    \bigcup M\sur\OccQuotCarrier{M}}{%
    \FormulaRefAuto{QuotProjSurjective}{1}}
\end{tabproof}

\FormulaThmDeltaK[Wert der Ununterscheidbarkeitsprojektion]{%
  a\in\bigcup M\vdash
  \OccQuotProj{M}(a)=\OccClass{M}{a}
}{OccQuotProjValue}{%
  \DeltaRow{Mengen}{M\dsep a}
  \DeltaRow{Relationsoperatoren}{\SameOccRel{M}}
  \DeltaRow{Funktionen}{\OccQuotProj{M}}
}
\begin{tabproof}
  \proofstep{1}{a\in\bigcup M}{\rA}
  \proofstep{}{\EqRel{\bigcup M,\SameOccRel{M}}}{%
    \FormulaRefAuto{SameOccRelEqRel}}
  \proofstep{1}{\OccQuotProj{M}(a)=\OccClass{M}{a}}{%
    \FormulaRefAuto{QuotProjValue}{2,1}}
\end{tabproof}

\FormulaThmDeltaK[Gleichheit von Ununterscheidbarkeitsprojektionswerten]{%
  a\in\bigcup M\dsep b\in\bigcup M\vdash
  \bigl(\OccQuotProj{M}(a)=\OccQuotProj{M}(b)\bigr)
  \leftrightarrow a\mathrel{\SameOccRel{M}}b
}{OccQuotProjEqualIffSameOcc}{%
  \DeltaRow{Mengen}{M\dsep a\dsep b}
  \DeltaRow{Relationsoperatoren}{\SameOccRel{M}}
  \DeltaRow{Funktionen}{\OccQuotProj{M}}
}
\begin{tabproofwide}
  \proofstepwidestar[1]{a\in\bigcup M}{\rA}
  \proofstepwidestar[2]{b\in\bigcup M}{\rA}
  \proofstepwidestar[]{\EqRel{\bigcup M,\SameOccRel{M}}}{%
    \FormulaRefAuto{SameOccRelEqRel}}
  \proofstepwidestar[1,2]{%
    \bigl(\OccQuotProj{M}(a)=\OccQuotProj{M}(b)\bigr)
    \leftrightarrow a\mathrel{\SameOccRel{M}}b}{%
    \FormulaRefAuto{QuotProjEqualIffRel}{3,1,2}}
\end{tabproofwide}

\subsection{Die induzierte Quotientenbildabbildung}

Die Abbildung \(\OccQuotMap{M}\) wendet die punktweise Projektion auf alle
Elemente eines Familienglieds an. Sie ist also die Einschränkung der von
\(\OccQuotProj{M}\) induzierten Potenzmengenabbildung auf \(M\).

Die folgenden kurzen Typisierungs- und Charakterisierungssätze dienen nur als
Schnittstelle für spätere Bände. Sie sind unmittelbare Spezialisierungen der
allgemeinen Sätze dieses Kapitels; neue inhaltliche Arbeit beginnt erst mit dem
Mitgliedschaftstransport.

\FormulaThmDeltaK[Typisierung der induzierten Ununterscheidbarkeits-Quotientenbildabbildung]{%
  \OccQuotMap{M}\colon M\to\powerset(\OccQuotCarrier{M})
}{OccQuotMapPowerFunction}{%
  \DeltaRow{Mengen}{M}
  \DeltaRow{Relationsoperatoren}{\SameOccRel{M}}
  \DeltaRow{Funktionen}{\OccQuotMap{M}}
}
\begin{tabproof}
  \proofstep{}{\EqRel{\bigcup M,\SameOccRel{M}}}{%
    \FormulaRefAuto{SameOccRelEqRel}}
  \proofstep{}{M\subseteq\powerset(\bigcup M)}{%
    \FormulaRefAuto{FamilySubsetPowerUnion}}
  \proofstep{}{\OccQuotMap{M}\colon
    M\to\powerset(\OccQuotCarrier{M})}{%
    \FormulaRefAuto{QuotFamMapFunction}{1,2}}
\end{tabproof}

\FormulaThmDeltaK[Charakterisierung der Quotientenbilder]{%
  X\in M\vdash
  \begin{aligned}[t]
    C\in\OccQuotMap{M}(X)
    \leftrightarrow{}&
    C\in\OccQuotCarrier{M}\land{}\\
    &\exists a\in X\,C=\OccClass{M}{a}
  \end{aligned}
}{OccQuotMapChar}{%
  \DeltaRow{Mengen}{M\dsep X\dsep C\dsep a}
  \DeltaRow{Relationsoperatoren}{\SameOccRel{M}}
  \DeltaRow{Funktionen}{\OccQuotMap{M}}
}
\begin{tabproof}
  \proofstep{1}{X\in M}{\rA}
  \proofstep{}{\EqRel{\bigcup M,\SameOccRel{M}}}{%
    \FormulaRefAuto{SameOccRelEqRel}}
  \proofstep{}{M\subseteq\powerset(\bigcup M)}{%
    \FormulaRefAuto{FamilySubsetPowerUnion}}
  \proofstep{1}{%
    C\in\OccQuotMap{M}(X)\leftrightarrow
    \bigl(C\in\OccQuotCarrier{M}\land
    \exists a\in X\,C=\OccClass{M}{a}\bigr)}{%
    \FormulaRefAuto{QuotFamMapChar}{2,3,1}}
\end{tabproof}

\FormulaThmDeltaK[Typisierung auf die Quotientenfamilie]{%
  \OccQuotMap{M}\colon M\to\OccQuotFam{M}
}{OccQuotMapFunction}{%
  \DeltaRow{Mengen}{M}
  \DeltaRow{Funktionen}{\OccQuotMap{M}}
}
\begin{tabproof}
  \proofstep{}{\OccQuotMap{M}\colon
    M\to\powerset(\OccQuotCarrier{M})}{%
    \FormulaRefAuto{OccQuotMapPowerFunction}}
  \proofstep{}{\OccQuotMap{M}\colon M\sur\OccQuotFam{M}}{%
    \FormulaRefAuto{CorestrictionToImageSurjective}{1}}
  \proofstep{}{\OccQuotMap{M}\colon M\to\OccQuotFam{M}}{%
    \FormulaRefAuto{Surjektive Funktion}{2}}
\end{tabproof}

\FormulaThmDeltaK[Quotientenbilder gehören zur Quotientenfamilie]{%
  X\in M\vdash \OccQuotMap{M}(X)\in\OccQuotFam{M}
}{OccQuotSetInFam}{%
  \DeltaRow{Mengen}{M\dsep X}
  \DeltaRow{Funktionen}{\OccQuotMap{M}}
}
\begin{tabproof}
  \proofstep{1}{X\in M}{\rA}
  \proofstep{}{\OccQuotMap{M}\colon M\to\OccQuotFam{M}}{%
    \FormulaRefAuto{OccQuotMapFunction}}
  \proofstep{1}{\OccQuotMap{M}(X)\in\OccQuotFam{M}}{%
    \FormulaRefAuto{F\colon A\to B\dsep x\in A\vdash F(x)\in B}{2,1}}
\end{tabproof}

\FormulaThmDeltaK[Mitglieder der Quotientenfamilie besitzen Urbilder]{%
  Y\in\OccQuotFam{M}\vdash
  \exists X\in M\,Y=\OccQuotMap{M}(X)
}{OccQuotFamPreimage}{%
  \DeltaRow{Mengen}{M\dsep X\dsep Y}
  \DeltaRow{Funktionen}{\OccQuotMap{M}}
}
\begin{tabproof}
  \proofstep{1}{Y\in\OccQuotFam{M}}{\rA}
  \proofstep{}{\OccQuotMap{M}\colon
    M\to\powerset(\OccQuotCarrier{M})}{%
    \FormulaRefAuto{OccQuotMapPowerFunction}}
  \proofstep{1}{\exists X\in M\,Y=\OccQuotMap{M}(X)}{%
    \FormulaRefAuto{F\colon A\to B\dsep y\in F[A]\vdash
    \exists x\in A\,y=F(x)}{2,1}}
\end{tabproof}

\FormulaThmDeltaK[Mitglieder der Quotientenfamilie liegen im Träger]{%
  Y\in\OccQuotFam{M}\vdash Y\subseteq\OccQuotCarrier{M}
}{OccQuotMemberSubsetCarrier}{%
  \DeltaRow{Mengen}{M\dsep Y}
  \DeltaRow{Funktionen}{\OccQuotMap{M}}
}
\begin{tabproof}
  \proofstep{1}{Y\in\OccQuotFam{M}}{\rA}
  \proofstep{}{\OccQuotMap{M}\colon
    M\to\powerset(\OccQuotCarrier{M})}{%
    \FormulaRefAuto{OccQuotMapPowerFunction}}
  \proofstep{}{\OccQuotMap{M}[M]\subseteq
    \powerset(\OccQuotCarrier{M})}{%
    \FormulaRefAuto{F\colon A\to B\vdash F[A]\subseteq B}{2}}
  \proofstep{1}{Y\in\powerset(\OccQuotCarrier{M})}{%
    \FormulaRefAuto{A\subseteq B,\,x\in A\vdash x\in B}{3,1}}
  \proofstep{1}{Y\subseteq\OccQuotCarrier{M}}{%
    \FormulaRefAuto{x\in\powerset(A)\eqvdash x\subseteq A}{4}}
\end{tabproof}

\paragraph{Struktureigenschaften}

Ab hier werden die für Mengenfamilien wesentlichen Folgen der Spezialisierung
hergeleitet: Mitgliedschaftstransport, Trennung der Familienglieder,
Vereinigungsverträglichkeit und Bijektivität.

\FormulaThmDeltaK[Mitgliedschaftstransport in den Quotienten]{%
  X\in M\dsep a\in\bigcup M
  \vdash
  a\in X\leftrightarrow
  \OccClass{M}{a}\in\OccQuotMap{M}(X)
}{OccQuotMembershipTransport}{%
  \DeltaRow{Mengen}{M\dsep X\dsep a\dsep b}
  \DeltaRow{Relationsoperatoren}{\SameOccRel{M}}
  \DeltaRow{Funktionen}{\OccQuotMap{M}}
}
\begin{tabproofsplitwide}
  \proofpartwideindR[Hinrichtung]{%
    X\in M\dsep a\in\bigcup M
    \vdash
    a\in X\rightarrow\OccClass{M}{a}\in\OccQuotMap{M}(X)
  }{%
    X\in M\dsep a\in\bigcup M
    \vdash
    a\in X\rightarrow\OccClass{M}{a}\in\OccQuotMap{M}(X)
  }[OccQuotMembershipTransportForward]
    \proofstepwidestar[1]{X\in M}{\rA}
    \proofstepwidestar[2]{a\in\bigcup M}{\rA}
    \proofstepwidestar[]{\EqRel{\bigcup M,\SameOccRel{M}}}{%
      \FormulaRefAuto{SameOccRelEqRel}}
    \proofstepwidestar[4]{a\in X}{\rA}
    \proofstepwidestar[2]{\OccClass{M}{a}\in\OccQuotCarrier{M}}{%
      \FormulaRefAuto{EqClassInQuotient}{3,2}}
    \proofstepwidestar[]{\OccClass{M}{a}=\OccClass{M}{a}}{\rII}
    \proofstepwidestar[4]{\exists b\in X\,
      \OccClass{M}{a}=\OccClass{M}{b}}{%
      \rEI{\rAI{4,6}}}
    \proofstepwidestar[2,4]{\OccClass{M}{a}\in\OccQuotCarrier{M}\land
      \exists b\in X\,\OccClass{M}{a}=\OccClass{M}{b}}{%
      \rAI{5,7}}
    \proofstepwidestar[1,2,4]{%
      \OccClass{M}{a}\in\OccQuotMap{M}(X)}{%
      \FormulaRefAuto{OccQuotMapChar}{1,8}}
    \proofstepwidestar[1,2]{%
      a\in X\rightarrow\OccClass{M}{a}\in\OccQuotMap{M}(X)}{%
      \rRI{4,9}}
  \closeproofpartwideindR

  \proofpartwideindR[Rückrichtung]{%
    X\in M\dsep a\in\bigcup M
    \vdash
    \OccClass{M}{a}\in\OccQuotMap{M}(X)\rightarrow a\in X
  }{%
    X\in M\dsep a\in\bigcup M
    \vdash
    \OccClass{M}{a}\in\OccQuotMap{M}(X)\rightarrow a\in X
  }[OccQuotMembershipTransportBackward]
    \proofstepwidestar[1]{X\in M}{\rA}
    \proofstepwidestar[2]{a\in\bigcup M}{\rA}
    \proofstepwidestar[]{\EqRel{\bigcup M,\SameOccRel{M}}}{%
      \FormulaRefAuto{SameOccRelEqRel}}
    \proofstepwidestar[4]{%
      \OccClass{M}{a}\in\OccQuotMap{M}(X)}{\rA}
    \proofstepwidestar[1,4]{\OccClass{M}{a}\in\OccQuotCarrier{M}\land
      \exists b\in X\,\OccClass{M}{a}=\OccClass{M}{b}}{%
      \FormulaRefAuto{OccQuotMapChar}{1,4}}
    \proofstepwidestar[1,4]{\exists b\in X\,
      \OccClass{M}{a}=\OccClass{M}{b}}{\rAEb{5}}
    \proofstepwidestar[7]{b\in X\land
      \OccClass{M}{a}=\OccClass{M}{b}}{\rA}
    \proofstepwidestar[7]{b\in X}{\rAEa{7}}
    \proofstepwidestar[1,7]{b\in\bigcup M}{%
      \FormulaRefAuto{B\in A\dsep x\in B\vdash x\in\bigcup A}{1,8}}
    \proofstepwidestar[7]{\OccClass{M}{a}=\OccClass{M}{b}}{\rAEb{7}}
    \proofstepwidestar[1,2,7]{a\mathrel{\SameOccRel{M}}b}{%
      \FormulaRefAuto{EqClassRelFromEqual}{3,2,9,10}}
    \proofstepwidestar[1,2,7]{%
      \forall Z\in M\,(a\in Z\leftrightarrow b\in Z)}{%
      \FormulaRefAuto{P\leftrightarrow Q, P\vdash Q}{%
        \FormulaRefAuto{SameOccRelChar}[thm]{2,9},11}}
    \proofstepwidestar[1,2,7]{a\in X\leftrightarrow b\in X}{%
      \rRE{1,\rUE{12}}}
    \proofstepwidestar[1,2,7]{a\in X}{%
      \FormulaRefAuto{P\leftrightarrow Q,\,Q\vdash P}{13,8}}
    \proofstepwidestar[1,2,4]{a\in X}{%
      \rEE{6,7,14}}
    \proofstepwidestar[1,2]{%
      \OccClass{M}{a}\in\OccQuotMap{M}(X)\rightarrow a\in X}{%
      \rRI{4,15}}
  \closeproofpartwideindR

  \proofpartwide{Schluss}
    \proofstepwidestar[1]{X\in M}{\rA}
    \proofstepwidestar[2]{a\in\bigcup M}{\rA}
    \proofstepwidestar[1,2]{%
      a\in X\rightarrow\OccClass{M}{a}\in\OccQuotMap{M}(X)}{%
      \FormulaRefAuto{OccQuotMembershipTransportForward}{1,2}}
    \proofstepwidestar[1,2]{%
      \OccClass{M}{a}\in\OccQuotMap{M}(X)\rightarrow a\in X}{%
      \FormulaRefAuto{OccQuotMembershipTransportBackward}{1,2}}
    \proofstepwidestar[1,2]{%
      a\in X\leftrightarrow
      \OccClass{M}{a}\in\OccQuotMap{M}(X)}{%
      \rLRI{3,4}}
  \closeproofpartwide
\end{tabproofsplitwide}

\FormulaThmDeltaK[Quotientenbilder trennen Familienglieder]{%
  X\in M\dsep Y\in M\dsep
  \OccQuotMap{M}(X)=\OccQuotMap{M}(Y)
  \vdash X=Y
}{OccQuotSetSeparatesMembers}{%
  \DeltaRow{Mengen}{M\dsep X\dsep Y\dsep z}
  \DeltaRow{Funktionen}{\OccQuotMap{M}}
}
\begin{tabproofsplitwide}
  \proofpartwideindR[Teilmengenrichtung \(X\subseteq Y\)]{%
    X\in M\dsep Y\in M\dsep
    \OccQuotMap{M}(X)=\OccQuotMap{M}(Y)
    \vdash X\subseteq Y
  }{%
    X\in M\dsep Y\in M\dsep
    \OccQuotMap{M}(X)=\OccQuotMap{M}(Y)
    \vdash X\subseteq Y
  }[OccQuotSetSeparatesMembersSubsetForward]
    \proofstepwidestar[1]{X\in M}{\rA}
    \proofstepwidestar[2]{Y\in M}{\rA}
    \proofstepwidestar[3]{\OccQuotMap{M}(X)=\OccQuotMap{M}(Y)}{\rA}
    \proofstepwidestar[4]{z\in X}{\rA}
    \proofstepwidestar[1,4]{z\in\bigcup M}{%
      \FormulaRefAuto{B\in A\dsep x\in B\vdash x\in\bigcup A}{1,4}}
    \proofstepwidestar[1,4]{%
      z\in X\leftrightarrow
      \OccClass{M}{z}\in\OccQuotMap{M}(X)}{%
      \FormulaRefAuto{OccQuotMembershipTransport}{1,5}}
    \proofstepwidestar[1,4]{%
      \OccClass{M}{z}\in\OccQuotMap{M}(X)}{%
      \FormulaRefAuto{P\leftrightarrow Q,\,P\vdash Q}{6,4}}
    \proofstepwidestar[1,3,4]{%
      \OccClass{M}{z}\in\OccQuotMap{M}(Y)}{\rIE{3,7}}
    \proofstepwidestar[1,2,4]{%
      z\in Y\leftrightarrow
      \OccClass{M}{z}\in\OccQuotMap{M}(Y)}{%
      \FormulaRefAuto{OccQuotMembershipTransport}{2,5}}
    \proofstepwidestar[1,2,3,4]{z\in Y}{%
      \FormulaRefAuto{P\leftrightarrow Q,\,Q\vdash P}{9,8}}
    \proofstepwidestar[1,2,3]{X\subseteq Y}{%
      \FormulaRefAuto{A\subseteq B\coloneq
      \forall z\,(z\in A\rightarrow z\in B)}{\rUI{\rRI{4,10}}}}
  \closeproofpartwideindR

  \proofpartwideindR[Teilmengenrichtung \(Y\subseteq X\)]{%
    X\in M\dsep Y\in M\dsep
    \OccQuotMap{M}(X)=\OccQuotMap{M}(Y)
    \vdash Y\subseteq X
  }{%
    X\in M\dsep Y\in M\dsep
    \OccQuotMap{M}(X)=\OccQuotMap{M}(Y)
    \vdash Y\subseteq X
  }[OccQuotSetSeparatesMembersSubsetBackward]
    \proofstepwidestar[1]{X\in M}{\rA}
    \proofstepwidestar[2]{Y\in M}{\rA}
    \proofstepwidestar[3]{\OccQuotMap{M}(X)=\OccQuotMap{M}(Y)}{\rA}
    \proofstepwidestar[3]{\OccQuotMap{M}(Y)=\OccQuotMap{M}(X)}{%
      \FormulaRefAuto{a=b\vdash b=a}{3}}
    \proofstepwidestar[1,2,3]{Y\subseteq X}{%
      \FormulaRefAuto{OccQuotSetSeparatesMembersSubsetForward}{2,1,4}}
  \closeproofpartwideindR

  \proofpartwide{Schluss}
    \proofstepwidestar[1]{X\in M}{\rA}
    \proofstepwidestar[2]{Y\in M}{\rA}
    \proofstepwidestar[3]{\OccQuotMap{M}(X)=\OccQuotMap{M}(Y)}{\rA}
    \proofstepwidestar[1,2,3]{X\subseteq Y}{%
      \FormulaRefAuto{OccQuotSetSeparatesMembersSubsetForward}{1,2,3}}
    \proofstepwidestar[1,2,3]{Y\subseteq X}{%
      \FormulaRefAuto{OccQuotSetSeparatesMembersSubsetBackward}{1,2,3}}
    \proofstepwidestar[1,2,3]{X=Y}{%
      \FormulaRefAuto{A\subseteq B, B\subseteq A\vdash A=B}{4,5}}
  \closeproofpartwide
\end{tabproofsplitwide}

\FormulaThmDeltaK[Vereinigungsverträglichkeit des Ununterscheidbarkeitsquotienten]{%
  X\in M\dsep Y\in M\dsep X\cup Y\in M
  \vdash
  \OccQuotMap{M}(X\cup Y)=
  \OccQuotMap{M}(X)\cup\OccQuotMap{M}(Y)
}{OccQuotPreservesUnion}{%
  \DeltaRow{Mengen}{M\dsep X\dsep Y}
  \DeltaRow{Funktionen}{\OccQuotMap{M}}
}
\begin{tabproof}
  \proofstep{1}{X\in M}{\rA}
  \proofstep{2}{Y\in M}{\rA}
  \proofstep{3}{X\cup Y\in M}{\rA}
  \proofstep{}{\EqRel{\bigcup M,\SameOccRel{M}}}{%
    \FormulaRefAuto{SameOccRelEqRel}}
  \proofstep{}{M\subseteq\powerset(\bigcup M)}{%
    \FormulaRefAuto{FamilySubsetPowerUnion}}
  \proofstep{1,2,3}{%
    \OccQuotMap{M}(X\cup Y)=
    \OccQuotMap{M}(X)\cup\OccQuotMap{M}(Y)}{%
    \FormulaRefAuto{QuotFamMapPreservesUnion}{4,5,1,2,3}}
\end{tabproof}

\FormulaThmDeltaK[Bijektivität der Ununterscheidbarkeitsquotientenabbildung]{%
  \OccQuotMap{M}\colon M\bij\OccQuotFam{M}
}{OccQuotMapBijective}{%
  \DeltaRow{Mengen}{M\dsep X\dsep Y}
  \DeltaRow{Funktionen}{\OccQuotMap{M}}
}
\begin{tabproof}
  \proofstep{}{\OccQuotMap{M}\colon M\to\OccQuotFam{M}}{%
    \FormulaRefAuto{OccQuotMapFunction}}
  \proofstep{2}{X\in M}{\rA}
  \proofstep{3}{Y\in M}{\rA}
  \proofstep{4}{\OccQuotMap{M}(X)=\OccQuotMap{M}(Y)}{\rA}
  \proofstep{2,3,4}{X=Y}{%
    \FormulaRefAuto{OccQuotSetSeparatesMembers}{2,3,4}}
  \proofstep{}{\OccQuotMap{M}\colon M\inj\OccQuotFam{M}}{%
    \FormulaRefAuto{Injektive Funktion}{1,5}}
  \proofstep{}{\OccQuotMap{M}\colon
    M\to\powerset(\OccQuotCarrier{M})}{%
    \FormulaRefAuto{OccQuotMapPowerFunction}}
  \proofstep{}{\OccQuotMap{M}\colon M\sur\OccQuotFam{M}}{%
    \FormulaRefAuto{CorestrictionToImageSurjective}{7}}
  \proofstep{}{\OccQuotMap{M}\colon M\bij\OccQuotFam{M}}{%
    \FormulaRefAuto{F\colon A\inj B\dsep F\colon A\sur B
    \vdash F\colon A\bij B}{6,8}}
\end{tabproof}

\FormulaThmDeltaK[Vereinigungsabschluss der Quotientenfamilie]{%
  \UnionClosedFam(M)\vdash\UnionClosedFam(\OccQuotFam{M})
}{OccQuotUnionClosed}{%
  \DeltaRow{Mengen}{M\dsep X\dsep Y\dsep U\dsep V}
  \DeltaRow{Funktionen}{\OccQuotMap{M}}
}
\begin{tabproofsplitwide}
  \proofpartwideindR[Transport fester Urbilder]{%
    \begin{aligned}[t]
      &\UnionClosedFam(M)\dsep
      X\in M\land U=\OccQuotMap{M}(X)\dsep{}\\[-2pt]
      &Y\in M\land V=\OccQuotMap{M}(Y)
      \vdash U\cup V\in\OccQuotFam{M}
    \end{aligned}
  }{%
    \UnionClosedFam(M)\dsep
    X\in M\land U=\OccQuotMap{M}(X)\dsep
    Y\in M\land V=\OccQuotMap{M}(Y)
    \vdash U\cup V\in\OccQuotFam{M}
  }[OccQuotUnionClosedWitnessTransport]
    \proofstepwidestar[1]{\UnionClosedFam(M)}{\rA}
    \proofstepwidestar[2]{X\in M\land U=\OccQuotMap{M}(X)}{\rA}
    \proofstepwidestar[3]{Y\in M\land V=\OccQuotMap{M}(Y)}{\rA}
    \proofstepwidestar[2]{X\in M}{\rAEa{2}}
    \proofstepwidestar[3]{Y\in M}{\rAEa{3}}
    \proofstepwidestar[1,2,3]{X\cup Y\in M}{%
      \FormulaRefAuto{Vereinigungsabgeschlossene Familie}{1,4,5}}
    \proofstepwidestar[2,3]{%
      U\cup V=\OccQuotMap{M}(X)\cup\OccQuotMap{M}(Y)}{%
      \rIE{\rAEb{2},\rAEb{3}}}
    \proofstepwidestar[1,2,3]{%
      \OccQuotMap{M}(X)\cup\OccQuotMap{M}(Y)=
      \OccQuotMap{M}(X\cup Y)}{%
      \FormulaRefAuto{a=b\vdash b=a}{%
        \FormulaRefAuto{OccQuotPreservesUnion}{4,5,6}}}
    \proofstepwidestar[1,2,3]{%
      U\cup V=\OccQuotMap{M}(X\cup Y)}{%
      \FormulaRefAuto{a = b,\, b = c \vdash a = c}{7,8}}
    \proofstepwidestar[1,2,3]{%
      \OccQuotMap{M}(X\cup Y)\in\OccQuotFam{M}}{%
      \FormulaRefAuto{OccQuotSetInFam}{6}}
    \proofstepwidestar[1,2,3]{U\cup V\in\OccQuotFam{M}}{%
      \rIE{9,10}}
  \closeproofpartwideindR

  \pagebreak[3]
  \proofpartwide{Schluss}
    \proofstepwidestar[1]{\UnionClosedFam(M)}{\rA}
    \proofstepwidestar[2]{U\in\OccQuotFam{M}}{\rA}
    \proofstepwidestar[3]{V\in\OccQuotFam{M}}{\rA}
    \proofstepwidestar[2]{\exists X\in M\,U=\OccQuotMap{M}(X)}{%
      \FormulaRefAuto{OccQuotFamPreimage}{2}}
    \proofstepwidestar[3]{\exists Y\in M\,V=\OccQuotMap{M}(Y)}{%
      \FormulaRefAuto{OccQuotFamPreimage}{3}}
    \proofstepwidestar[6]{X\in M\land U=\OccQuotMap{M}(X)}{\rA}
    \proofstepwidestar[7]{Y\in M\land V=\OccQuotMap{M}(Y)}{\rA}
    \proofstepwidestar[1,6,7]{U\cup V\in\OccQuotFam{M}}{%
      \FormulaRefAuto{OccQuotUnionClosedWitnessTransport}{1,6,7}}
    \proofstepwidestar[1,6]{U\cup V\in\OccQuotFam{M}}{%
      \rEE{5,7,8}}
    \proofstepwidestar[1]{U\cup V\in\OccQuotFam{M}}{%
      \rEE{4,6,9}}
    \proofstepwidestar[1]{\UnionClosedFam(\OccQuotFam{M})}{%
      \FormulaRefAuto{Vereinigungsabgeschlossene Familie}{%
        \rUI{\rRI{2,\rRI{3,10}}}}}
  \closeproofpartwide
\end{tabproofsplitwide}

\FormulaThmDeltaK[Vereinigung der Quotientenfamilie]{%
  \bigcup\OccQuotFam{M}=\OccQuotCarrier{M}
}{OccQuotFamUnionCarrier}{%
  \DeltaRow{Mengen}{M\dsep X\dsep Y\dsep C\dsep a}
  \DeltaRow{Relationsoperatoren}{\SameOccRel{M}}
  \DeltaRow{Funktionen}{\OccQuotMap{M}}
}
\begin{tabproofsplitwide}
  \proofpartwideindR[Die Vereinigung liegt im Quotiententräger]{%
    \bigcup\OccQuotFam{M}\subseteq\OccQuotCarrier{M}
  }{%
    \bigcup\OccQuotFam{M}\subseteq\OccQuotCarrier{M}
  }[OccQuotFamUnionCarrierSubsetForward]
    \proofstepwidestar[1]{C\in\bigcup\OccQuotFam{M}}{\rA}
    \proofstepwidestar[1]{\exists Y\in\OccQuotFam{M}\,C\in Y}{%
      \FormulaRefAuto{x\in\bigcup A\eqvdash
      \exists B\in A\,(x\in B)}{1}}
    \proofstepwidestar[3]{Y\in\OccQuotFam{M}\land C\in Y}{\rA}
    \proofstepwidestar[3]{Y\in\OccQuotFam{M}}{\rAEa{3}}
    \proofstepwidestar[3]{C\in Y}{\rAEb{3}}
    \proofstepwidestar[3]{\exists X\in M\,Y=\OccQuotMap{M}(X)}{%
      \FormulaRefAuto{OccQuotFamPreimage}{4}}
    \proofstepwidestar[7]{X\in M\land Y=\OccQuotMap{M}(X)}{\rA}
    \proofstepwidestar[7]{X\in M}{\rAEa{7}}
    \proofstepwidestar[3,7]{C\in\OccQuotMap{M}(X)}{%
      \rIE{\rAEb{7},5}}
    \proofstepwidestar[3,7]{C\in\OccQuotCarrier{M}\land
      \exists a\in X\,C=\OccClass{M}{a}}{%
      \FormulaRefAuto{OccQuotMapChar}{8,9}}
    \proofstepwidestar[3,7]{C\in\OccQuotCarrier{M}}{\rAEa{10}}
    \proofstepwidestar[3]{C\in\OccQuotCarrier{M}}{%
      \rEE{6,7,11}}
    \proofstepwidestar[1]{C\in\OccQuotCarrier{M}}{%
      \rEE{2,3,12}}
    \proofstepwidestar[]{%
      \bigcup\OccQuotFam{M}\subseteq\OccQuotCarrier{M}}{%
      \FormulaRefAuto{A\subseteq B\coloneq
      \forall C\,(C\in A\rightarrow C\in B)}{\rUI{\rRI{1,13}}}}
  \closeproofpartwideindR

  \proofpartwideindR[Jede Quotientenklasse liegt in der Vereinigung]{%
    \OccQuotCarrier{M}\subseteq\bigcup\OccQuotFam{M}
  }{%
    \OccQuotCarrier{M}\subseteq\bigcup\OccQuotFam{M}
  }[OccQuotFamUnionCarrierSubsetBackward]
    \proofstepwidestar[1]{C\in\OccQuotCarrier{M}}{\rA}
    \proofstepwidestar[]{\EqRel{\bigcup M,\SameOccRel{M}}}{%
      \FormulaRefAuto{SameOccRelEqRel}}
    \proofstepwidestar[1]{\exists a\in\bigcup M\,
      C=\OccClass{M}{a}}{%
      \FormulaRefAuto{QuotientRepresentative}{2,1}}
    \proofstepwidestar[4]{a\in\bigcup M\land C=\OccClass{M}{a}}{\rA}
    \proofstepwidestar[4]{a\in\bigcup M}{\rAEa{4}}
    \proofstepwidestar[4]{C=\OccClass{M}{a}}{\rAEb{4}}
    \proofstepwidestar[4]{\exists X\in M\,a\in X}{%
      \FormulaRefAuto{x\in\bigcup A\eqvdash
      \exists B\in A\,(x\in B)}{5}}
    \proofstepwidestar[8]{X\in M\land a\in X}{\rA}
    \proofstepwidestar[8]{X\in M}{\rAEa{8}}
    \proofstepwidestar[8]{a\in X}{\rAEb{8}}
    \proofstepwidestar[4,8]{%
      a\in X\leftrightarrow
      \OccClass{M}{a}\in\OccQuotMap{M}(X)}{%
      \FormulaRefAuto{OccQuotMembershipTransport}{9,5}}
    \proofstepwidestar[4,8]{%
      \OccClass{M}{a}\in\OccQuotMap{M}(X)}{%
      \FormulaRefAuto{P\leftrightarrow Q,\,P\vdash Q}{11,10}}
    \proofstepwidestar[4,8]{C\in\OccQuotMap{M}(X)}{%
      \rIE{6,12}}
    \proofstepwidestar[8]{\OccQuotMap{M}(X)\in\OccQuotFam{M}}{%
      \FormulaRefAuto{OccQuotSetInFam}{9}}
    \proofstepwidestar[4,8]{C\in\bigcup\OccQuotFam{M}}{%
      \FormulaRefAuto{x\in\bigcup A\eqvdash
      \exists B\in A\,(x\in B)}{\rEI{\rAI{14,13}}}}
    \proofstepwidestar[4]{C\in\bigcup\OccQuotFam{M}}{%
      \rEE{7,8,15}}
    \proofstepwidestar[1]{C\in\bigcup\OccQuotFam{M}}{%
      \rEE{3,4,16}}
    \proofstepwidestar[]{%
      \OccQuotCarrier{M}\subseteq\bigcup\OccQuotFam{M}}{%
      \FormulaRefAuto{A\subseteq B\coloneq
      \forall C\,(C\in A\rightarrow C\in B)}{\rUI{\rRI{1,17}}}}
  \closeproofpartwideindR

  \proofpartwide{Schluss}
    \proofstepwidestar[]{%
      \bigcup\OccQuotFam{M}\subseteq\OccQuotCarrier{M}}{%
      \FormulaRefAuto{OccQuotFamUnionCarrierSubsetForward}}
    \proofstepwidestar[]{%
      \OccQuotCarrier{M}\subseteq\bigcup\OccQuotFam{M}}{%
      \FormulaRefAuto{OccQuotFamUnionCarrierSubsetBackward}}
    \proofstepwidestar[]{%
      \bigcup\OccQuotFam{M}=\OccQuotCarrier{M}}{%
      \FormulaRefAuto{A\subseteq B, B\subseteq A\vdash A=B}{1,2}}
  \closeproofpartwide
\end{tabproofsplitwide}

\FormulaThmDeltaK[Elemente der vereinigten Quotientenfamilie haben Repräsentanten]{%
  C\in\bigcup\OccQuotFam{M}
  \vdash
  \exists a\in\bigcup M\,C=\OccClass{M}{a}
}{OccQuotUnionRepresentative}{%
  \DeltaRow{Mengen}{M\dsep C\dsep a}
  \DeltaRow{Relationsoperatoren}{\SameOccRel{M}}
}
\begin{tabproof}
  \proofstep{1}{C\in\bigcup\OccQuotFam{M}}{\rA}
  \proofstep{}{\bigcup\OccQuotFam{M}=\OccQuotCarrier{M}}{%
    \FormulaRefAuto{OccQuotFamUnionCarrier}}
  \proofstep{1}{C\in\OccQuotCarrier{M}}{\rIE{2,1}}
  \proofstep{}{\EqRel{\bigcup M,\SameOccRel{M}}}{%
    \FormulaRefAuto{SameOccRelEqRel}}
  \proofstep{1}{\exists a\in\bigcup M\,C=\OccClass{M}{a}}{%
    \FormulaRefAuto{QuotientRepresentative}{4,3}}
\end{tabproof}

\FormulaThmDeltaK[Nichtleerheit der Quotientenfamilie]{%
  M\neq\varnothing\vdash\OccQuotFam{M}\neq\varnothing
}{OccQuotFamNonempty}{%
  \DeltaRow{Mengen}{M\dsep X}
  \DeltaRow{Funktionen}{\OccQuotMap{M}}
}
\begin{tabproof}
  \proofstep{1}{M\neq\varnothing}{\rA}
  \proofstep{1}{\exists X\in M}{%
    \FormulaRefAuto{B\neq\varnothing\vdash\exists x\in B}{1}}
  \proofstep{3}{X\in M}{\rA}
  \proofstep{3}{\OccQuotMap{M}(X)\in\OccQuotFam{M}}{%
    \FormulaRefAuto{OccQuotSetInFam}{3}}
  \proofstep{3}{\OccQuotFam{M}\neq\varnothing}{%
    \FormulaRefAuto{a\in A\vdash A\neq\varnothing}{4}}
  \proofstep{1}{\OccQuotFam{M}\neq\varnothing}{\rEE{2,3,5}}
\end{tabproof}

\FormulaThmDeltaK[Nichtleerheit der Vereinigung der Quotientenfamilie]{%
  \bigcup M\neq\varnothing
  \vdash
  \bigcup\OccQuotFam{M}\neq\varnothing
}{OccQuotFamCarrierNonempty}{%
  \DeltaRow{Mengen}{M\dsep a}
  \DeltaRow{Relationsoperatoren}{\SameOccRel{M}}
}
\begin{tabproofsplit}
  \proofpartindR[Eine Klasse als Element der Vereinigung]{%
    a\in\bigcup M
    \vdash \OccClass{M}{a}\in\bigcup\OccQuotFam{M}%
  }{%
    a\in\bigcup M
    \vdash \OccClass{M}{a}\in\bigcup\OccQuotFam{M}%
  }[OccQuotFamCarrierElement]
    \proofstep{1}{a\in\bigcup M}{\rA}
    \proofstep{}{\EqRel{\bigcup M,\SameOccRel{M}}}{%
      \FormulaRefAuto{SameOccRelEqRel}}
    \proofstep{1}{\OccClass{M}{a}\in\OccQuotCarrier{M}}{%
      \FormulaRefAuto{EqClassInQuotient}{2,1}}
    \proofstep{}{\bigcup\OccQuotFam{M}=\OccQuotCarrier{M}}{%
      \FormulaRefAuto{OccQuotFamUnionCarrier}}
    \proofstep{}{\OccQuotCarrier{M}=\bigcup\OccQuotFam{M}}{%
      \FormulaRefAuto{a=b\vdash b=a}{4}}
    \proofstep{1}{\OccClass{M}{a}\in\bigcup\OccQuotFam{M}}{%
      \rIE{5,3}}
  \closeproofpartindR

  \proofpartindR[Existenz eines Elements]{%
    \bigcup M\neq\varnothing
    \vdash \bigcup\OccQuotFam{M}\neq\varnothing%
  }{%
    \bigcup M\neq\varnothing
    \vdash \bigcup\OccQuotFam{M}\neq\varnothing%
  }[OccQuotFamCarrierNonemptyConclusion]
    \proofstep{1}{\bigcup M\neq\varnothing}{\rA}
    \proofstep{1}{\exists a\in\bigcup M}{%
      \FormulaRefAuto{B\neq\varnothing\vdash\exists x\in B}{1}}
    \proofstep{3}{a\in\bigcup M}{\rA}
    \proofstep{3}{\OccClass{M}{a}\in\bigcup\OccQuotFam{M}}{%
      \FormulaRefAuto{OccQuotFamCarrierElement}{3}}
    \proofstep{3}{\bigcup\OccQuotFam{M}\neq\varnothing}{%
      \FormulaRefAuto{a\in A\vdash A\neq\varnothing}{4}}
    \proofstep{1}{\bigcup\OccQuotFam{M}\neq\varnothing}{%
      \rEE{2,3,5}}
  \closeproofpartindR
\end{tabproofsplit}

\paragraph{Die Spur trennt Klassen}

Die allgemeine Spurabbildung zeichnet jede Quotientenklasse durch die
Familienglieder aus, in denen sie vorkommt. Für den
Ununterscheidbarkeitsquotienten ist diese Kennzeichnung injektiv.

\FormulaThmDeltaK[Die allgemeine Spur trennt Ununterscheidbarkeitsklassen]{%
  \TraceMap{\OccQuotMap{M}}\colon
  \OccQuotCarrier{M}\inj\powerset(M)
}{OccQuotTraceMapInjective}{%
  \DeltaRow{Mengen}{M\dsep C\dsep D\dsep X\dsep a\dsep b}
  \DeltaRow{Relationsoperatoren}{\SameOccRel{M}}
  \DeltaRow{Funktionen}{\OccQuotMap{M}\dsep
    \TraceMap{\OccQuotMap{M}}}
}
\begin{tabproofsplitwide}
  \proofpartwideindR[Punktweiser Transport]{%
    \begin{aligned}[t]
      &C\in\OccQuotCarrier{M}\dsep D\in\OccQuotCarrier{M}\dsep
      \TraceMap{\OccQuotMap{M}}(C)=
      \TraceMap{\OccQuotMap{M}}(D)\dsep{}\\[-2pt]
      &a\in\bigcup M\dsep C=\OccClass{M}{a}\dsep
      b\in\bigcup M\dsep D=\OccClass{M}{b}\dsep X\in M\\
      &\vdash a\in X\leftrightarrow b\in X
    \end{aligned}
  }{%
    C\in\OccQuotCarrier{M}\dsep D\in\OccQuotCarrier{M}\dsep
    \TraceMap{\OccQuotMap{M}}(C)=
    \TraceMap{\OccQuotMap{M}}(D)\dsep
    a\in\bigcup M\dsep C=\OccClass{M}{a}\dsep
    b\in\bigcup M\dsep D=\OccClass{M}{b}\dsep X\in M
    \vdash a\in X\leftrightarrow b\in X
  }[OccQuotTraceRepresentativesMembership]
    \proofstepwidestar[]{\OccQuotMap{M}\colon
      M\to\powerset(\OccQuotCarrier{M})}{%
      \FormulaRefAuto{OccQuotMapPowerFunction}}
    \proofstepwidestar[2]{C\in\OccQuotCarrier{M}}{\rA}
    \proofstepwidestar[3]{D\in\OccQuotCarrier{M}}{\rA}
    \proofstepwidestar[4]{%
      \TraceMap{\OccQuotMap{M}}(C)=
      \TraceMap{\OccQuotMap{M}}(D)}{\rA}
    \proofstepwidestar[5]{a\in\bigcup M}{\rA}
    \proofstepwidestar[6]{C=\OccClass{M}{a}}{\rA}
    \proofstepwidestar[7]{b\in\bigcup M}{\rA}
    \proofstepwidestar[8]{D=\OccClass{M}{b}}{\rA}
    \proofstepwidestar[9]{X\in M}{\rA}
    \proofstepwidestar[1,2,3,4,9]{%
      C\in\OccQuotMap{M}(X)\leftrightarrow
      D\in\OccQuotMap{M}(X)}{%
      \FormulaRefAuto{TraceMapEqualityMembership}{%
        1,2,3,4,9}}
    \proofstepwidestar[5,9]{%
      a\in X\leftrightarrow
      \OccClass{M}{a}\in\OccQuotMap{M}(X)}{%
      \FormulaRefAuto{OccQuotMembershipTransport}{9,5}}
    \proofstepwidestar[5,6,9]{%
      a\in X\leftrightarrow C\in\OccQuotMap{M}(X)}{%
      \rIE{6,11}}
    \proofstepwidestar[7,9]{%
      b\in X\leftrightarrow
      \OccClass{M}{b}\in\OccQuotMap{M}(X)}{%
      \FormulaRefAuto{OccQuotMembershipTransport}{9,7}}
    \proofstepwidestar[7,8,9]{%
      b\in X\leftrightarrow D\in\OccQuotMap{M}(X)}{%
      \rIE{8,13}}
    \proofstepwidestar[7,8,9]{%
      D\in\OccQuotMap{M}(X)\leftrightarrow b\in X}{%
      \FormulaRefAuto{P\leftrightarrow Q\vdash Q\leftrightarrow P}{14}}
    \proofstepwidestar[1,2,3,4,5,6,9]{%
      a\in X\leftrightarrow D\in\OccQuotMap{M}(X)}{%
      \FormulaRefAuto{P\leftrightarrow Q\dsep Q\leftrightarrow R
      \vdash P\leftrightarrow R}{12,10}}
    \proofstepwidestar[1,2,3,4,5,6,7,8,9]{%
      a\in X\leftrightarrow b\in X}{%
      \FormulaRefAuto{P\leftrightarrow Q\dsep Q\leftrightarrow R
      \vdash P\leftrightarrow R}{16,15}}
  \closeproofpartwideindR

  \proofpartwideindR[Ununterscheidbarkeit der Repräsentanten]{%
    \begin{aligned}[t]
      &C,D\in\OccQuotCarrier{M}\dsep{}\\[-2pt]
      &\TraceMap{\OccQuotMap{M}}(C)=
      \TraceMap{\OccQuotMap{M}}(D)\dsep{}\\[-2pt]
      &a\in\bigcup M\dsep b\in\bigcup M\dsep{}\\[-2pt]
      &C=\OccClass{M}{a}\dsep D=\OccClass{M}{b}\\
      &\vdash a\mathrel{\SameOccRel{M}}b
    \end{aligned}
  }{%
    \begin{aligned}[t]
      &C,D\in\OccQuotCarrier{M}\dsep
        \TraceMap{\OccQuotMap{M}}(C)=
        \TraceMap{\OccQuotMap{M}}(D)\dsep{}\\[-2pt]
      &a,b\in\bigcup M\dsep C=\OccClass{M}{a}\dsep
        D=\OccClass{M}{b}
        \vdash a\mathrel{\SameOccRel{M}}b
    \end{aligned}
  }[OccQuotTraceRepresentativesRelated]
    \proofstepwidestar[1]{C\in\OccQuotCarrier{M}}{\rA}
    \proofstepwidestar[2]{D\in\OccQuotCarrier{M}}{\rA}
    \proofstepwidestar[3]{%
      \TraceMap{\OccQuotMap{M}}(C)=
      \TraceMap{\OccQuotMap{M}}(D)}{\rA}
    \proofstepwidestar[4]{a\in\bigcup M}{\rA}
    \proofstepwidestar[5]{C=\OccClass{M}{a}}{\rA}
    \proofstepwidestar[6]{b\in\bigcup M}{\rA}
    \proofstepwidestar[7]{D=\OccClass{M}{b}}{\rA}
    \proofstepwidestar[8]{X\in M}{\rA}
    \proofstepwidestar[1,2,3,4,5,6,7,8]{%
      a\in X\leftrightarrow b\in X}{%
      \FormulaRefAuto{OccQuotTraceRepresentativesMembership}{%
        1,2,3,4,5,6,7,8}}
    \proofstepwidestar[1,2,3,4,5,6,7]{%
      \forall X\in M\,(a\in X\leftrightarrow b\in X)}{%
      \rUI{\rRI{8,9}}}
    \proofstepwidestar[1,2,3,4,5,6,7]{%
      \begin{aligned}[t]
        &a\in\bigcup M\land b\in\bigcup M\land{}\\[-2pt]
        &\forall X\in M\,(a\in X\leftrightarrow b\in X)
      \end{aligned}}{%
      \rAI{4,\rAI{6,10}}}
    \proofstepwidestar[1,2,3,4,5,6,7]{a\mathrel{\SameOccRel{M}}b}{%
      \FormulaRefAuto{P\leftrightarrow Q, Q\vdash P}{%
        \FormulaRefAuto{SameOccRelChar}[thm]{4,6},10}}
  \closeproofpartwideindR

  \proofpartwideindR[Klassengleichheit]{%
    \begin{aligned}[t]
      &C,D\in\OccQuotCarrier{M}\dsep{}\\[-2pt]
      &\TraceMap{\OccQuotMap{M}}(C)=
      \TraceMap{\OccQuotMap{M}}(D)\\
      &\vdash C=D
    \end{aligned}
  }{%
    \begin{aligned}[t]
      &C,D\in\OccQuotCarrier{M}\dsep{}\\[-2pt]
      &\TraceMap{\OccQuotMap{M}}(C)=
        \TraceMap{\OccQuotMap{M}}(D)\vdash{}\\[-2pt]
      &C=D
    \end{aligned}
  }[OccQuotTraceMapEquality]
    \proofstepwidestar[1]{C\in\OccQuotCarrier{M}}{\rA}
    \proofstepwidestar[2]{D\in\OccQuotCarrier{M}}{\rA}
    \proofstepwidestar[3]{%
      \TraceMap{\OccQuotMap{M}}(C)=
      \TraceMap{\OccQuotMap{M}}(D)}{\rA}
    \proofstepwidestar[]{\EqRel{\bigcup M,\SameOccRel{M}}}{%
      \FormulaRefAuto{SameOccRelEqRel}}
    \proofstepwidestar[1]{\exists a\in\bigcup M\,
      C=\OccClass{M}{a}}{%
      \FormulaRefAuto{QuotientRepresentative}{4,1}}
    \proofstepwidestar[2]{\exists b\in\bigcup M\,
      D=\OccClass{M}{b}}{%
      \FormulaRefAuto{QuotientRepresentative}{4,2}}
    \proofstepwidestar[7]{a\in\bigcup M\land
      C=\OccClass{M}{a}}{\rA}
    \proofstepwidestar[8]{b\in\bigcup M\land
      D=\OccClass{M}{b}}{\rA}
    \proofstepwidestar[7]{a\in\bigcup M}{\rAEa{7}}
    \proofstepwidestar[7]{C=\OccClass{M}{a}}{\rAEb{7}}
    \proofstepwidestar[8]{b\in\bigcup M}{\rAEa{8}}
    \proofstepwidestar[8]{D=\OccClass{M}{b}}{\rAEb{8}}
    \proofstepwidestar[1,2,3,7,8]{a\mathrel{\SameOccRel{M}}b}{%
      \FormulaRefAuto{OccQuotTraceRepresentativesRelated}{%
        1,2,3,9,10,11,12}}
    \proofstepwidestar[1,2,3,7,8]{%
      \OccClass{M}{a}=\OccClass{M}{b}}{%
      \FormulaRefAuto{EqClassEqualFromRel}{%
        4,9,11,13}}
    \proofstepwidestar[1,2,3,7,8]{C=\OccClass{M}{b}}{%
      \FormulaRefAuto{a = b,\, b = c \vdash a = c}{10,14}}
    \proofstepwidestar[1,2,3,7,8]{C=D}{%
      \FormulaRefAuto{a = b,\, c = b \vdash a = c}{15,12}}
    \proofstepwidestar[1,2,3,7]{C=D}{\rEE{6,8,16}}
    \proofstepwidestar[1,2,3]{C=D}{\rEE{5,7,17}}
  \closeproofpartwideindR

  \pagebreak[3]
  \proofpartwide{Schluss}
    \proofstepwidestar[]{\OccQuotMap{M}\colon
      M\to\powerset(\OccQuotCarrier{M})}{%
      \FormulaRefAuto{OccQuotMapPowerFunction}}
    \proofstepwidestar[]{\TraceMap{\OccQuotMap{M}}\colon
      \OccQuotCarrier{M}\to\powerset(M)}{%
      \FormulaRefAuto{TraceMapFunction}{1}}
    \proofstepwidestar[3]{C\in\OccQuotCarrier{M}}{\rA}
    \proofstepwidestar[4]{D\in\OccQuotCarrier{M}}{\rA}
    \proofstepwidestar[5]{%
      \TraceMap{\OccQuotMap{M}}(C)=
      \TraceMap{\OccQuotMap{M}}(D)}{\rA}
    \proofstepwidestar[3,4,5]{C=D}{%
      \FormulaRefAuto{OccQuotTraceMapEquality}{3,4,5}}
    \proofstepwidestar[]{\TraceMap{\OccQuotMap{M}}\colon
      \OccQuotCarrier{M}\inj\powerset(M)}{%
      \FormulaRefAuto{Injektive Funktion}{2,6}}
  \closeproofpartwide
\end{tabproofsplitwide}

\paragraph{Beispiel.}
Für \(M=\{\varnothing,\{a,b\}\}\) treten \(a\) und \(b\) in genau denselben
Familiengliedern auf. Daher bilden sie dieselbe Klasse
\(C=\OccClass{M}{a}=\OccClass{M}{b}\). Es gilt
\[
  \OccQuotMap{M}(\varnothing)=\varnothing,\qquad
  \OccQuotMap{M}(\{a,b\})=\{C\},
  \qquad
  \OccQuotFam{M}=\{\varnothing,\{C\}\}.
\]


\end{document}
